\documentclass[10pt,a4paper,twoside,openany,parskip=half]{scrbook}

\usepackage[utf8]{inputenc}
\usepackage[T1]{fontenc}
\usepackage{amsmath,amsthm,mathtools}
\usepackage[lining,semibold]{libertine}
\usepackage[scaled=0.85]{beramono}
\let\openbox\relax
\let\Bbbk\relax
\usepackage[libertine,vvarbb]{newtxmath}
\let\Bbbk\relax
\usepackage{amssymb}
\usepackage[polish]{babel}
\usepackage{microtype}
\AtBeginDocument{\shorthandoff{"}}
\usepackage{tikz}
\usepackage{tikz-cd}
\usepackage{enumitem}
\setlist{itemsep=2pt,parsep=2pt,topsep=4pt}
\usepackage{xcolor}
\usepackage[most]{tcolorbox}
\usepackage[a4paper,inner=28mm,outer=24mm,top=24mm,bottom=26mm,headsep=8mm,footskip=14mm]{geometry}
\usepackage{scrlayer-scrpage}
\pagestyle{scrheadings}
\automark[section]{chapter}
\clearpairofpagestyles
\ihead{\textsc{\leftmark}}
\ohead{\textsc{\rightmark}}
\ofoot[\pagemark]{\pagemark}

\usepackage[colorlinks=true,linkcolor=teal!55!black,urlcolor=teal!55!black,citecolor=teal!55!black,bookmarksnumbered=true,bookmarksopen=true,pdftitle={Teoria kategorii},pdfauthor={Tomasz Brengos}]{hyperref}

% --- Palette ---
\definecolor{defaccent}{HTML}{1F4E79}
\definecolor{thmaccent}{HTML}{6A1B9A}
\definecolor{exaccent}{HTML}{2E7D32}
\definecolor{remaccent}{HTML}{546E7A}
\definecolor{zadaccent}{HTML}{D84315}
\definecolor{titleaccent}{HTML}{1F4E79}

% --- Chapter & section appearance ---
\setkomafont{disposition}{\normalfont\sffamily\bfseries}
\setkomafont{chapter}{\Huge\sffamily\bfseries\color{titleaccent}}
\setkomafont{chapterprefix}{\Large\sffamily\mdseries\color{remaccent}}
\setkomafont{section}{\Large\sffamily\bfseries\color{titleaccent!90!black}}
\setkomafont{subsection}{\large\sffamily\bfseries}
\renewcommand*{\chapterformat}{\textcolor{remaccent}{\thechapter}\enskip}
\renewcommand*{\raggedchapter}{\raggedright}

% --- Theorem environments ---
\theoremstyle{definition}
\newtheorem{definicja}{Definicja}[chapter]
\newtheorem{przyklad}[definicja]{Przykład}
\newtheorem{uwaga}[definicja]{Uwaga}
\newtheorem{fakt}[definicja]{Fakt}
\newtheorem{zadanie}{Zadanie}[section]
\theoremstyle{plain}
\newtheorem{twierdzenie}[definicja]{Twierdzenie}
\newtheorem{lemat}[definicja]{Lemat}
\newtheorem{wniosek}[definicja]{Wniosek}
\newtheorem{stwierdzenie}[definicja]{Stwierdzenie}

% Common box style: left accent bar, no surrounding frame, breakable
\tcolorboxenvironment{definicja}{
  blanker, breakable,
  before skip=10pt, after skip=10pt,
  borderline west={2.5pt}{0pt}{defaccent},
  left=10pt, top=2pt, bottom=2pt
}
\tcolorboxenvironment{twierdzenie}{
  blanker, breakable,
  before skip=10pt, after skip=10pt,
  borderline west={2.5pt}{0pt}{thmaccent},
  left=10pt, top=2pt, bottom=2pt
}
\tcolorboxenvironment{lemat}{
  blanker, breakable,
  before skip=10pt, after skip=10pt,
  borderline west={2.5pt}{0pt}{thmaccent},
  left=10pt, top=2pt, bottom=2pt
}
\tcolorboxenvironment{wniosek}{
  blanker, breakable,
  before skip=10pt, after skip=10pt,
  borderline west={2.5pt}{0pt}{thmaccent},
  left=10pt, top=2pt, bottom=2pt
}
\tcolorboxenvironment{stwierdzenie}{
  blanker, breakable,
  before skip=10pt, after skip=10pt,
  borderline west={2.5pt}{0pt}{thmaccent},
  left=10pt, top=2pt, bottom=2pt
}
\tcolorboxenvironment{przyklad}{
  blanker, breakable,
  before skip=10pt, after skip=10pt,
  borderline west={2.5pt}{0pt}{exaccent},
  left=10pt, top=2pt, bottom=2pt
}
\tcolorboxenvironment{uwaga}{
  blanker, breakable,
  before skip=8pt, after skip=8pt,
  borderline west={1.8pt}{0pt}{remaccent},
  left=8pt, top=2pt, bottom=2pt
}
\tcolorboxenvironment{fakt}{
  blanker, breakable,
  before skip=8pt, after skip=8pt,
  borderline west={1.8pt}{0pt}{remaccent},
  left=8pt, top=2pt, bottom=2pt
}
\tcolorboxenvironment{zadanie}{
  blanker, breakable,
  before skip=10pt, after skip=10pt,
  borderline west={2.5pt}{0pt}{zadaccent},
  left=10pt, top=2pt, bottom=2pt
}

% --- Math macros ---
\newcommand{\catname}[1]{\mathbf{#1}}
\newcommand{\Set}{\catname{Set}}
\newcommand{\Setf}{\catname{Set}_{\mathrm{fin}}}
\newcommand{\Grp}{\catname{Grp}}
\newcommand{\Ab}{\catname{Ab}}
\let\Top\relax
\newcommand{\Top}{\catname{Top}}
\newcommand{\Vect}{\catname{Vect}}
\newcommand{\Mon}{\catname{Mon}}
\newcommand{\Ring}{\catname{Ring}}
\newcommand{\Pos}{\catname{Pos}}
\newcommand{\Rel}{\catname{Rel}}
\newcommand{\Par}{\catname{Par}}
\newcommand{\Err}{\catname{Err}}
\newcommand{\Trace}{\catname{Trace}}
\newcommand{\Cat}{\mathcal{C}}
\let\C\undefined
\newcommand{\C}{\mathcal{C}}
\newcommand{\Dcat}{\mathcal{D}}
\newcommand{\D}{\mathcal{D}}
\newcommand{\Ecat}{\mathcal{E}}
\newcommand{\Jcat}{\mathcal{J}}
\newcommand{\Fcat}{\mathcal{F}}
\newcommand{\Ob}{\mathrm{Ob}}
\newcommand{\Mor}{\mathrm{Mor}}
\newcommand{\dom}{\mathrm{dom}}
\newcommand{\cod}{\mathrm{cod}}
\newcommand{\id}{\mathrm{id}}
\newcommand{\Id}{\mathrm{Id}}
\newcommand{\op}{\mathrm{op}}
\newcommand{\IO}{\mathrm{IO}}
\newcommand{\State}{\mathrm{State}}
\newcommand{\Kl}{\mathrm{Kl}}
\DeclareMathOperator{\Hom}{Hom}
\DeclareMathOperator{\Ker}{Ker}
\DeclareMathOperator{\Img}{Im}

% --- Front matter metadata ---
\titlehead{\centering\sffamily\small Politechnika Warszawska}
\subject{\sffamily Notatki z wykładu}
\title{\sffamily\bfseries Teoria kategorii}
\author{\sffamily Tomasz Brengos}
\date{}

\begin{document}
\frontmatter
\maketitle
\tableofcontents

\mainmatter

\chapter{Pojęcie kategorii. Przykłady i funktory}\label{ch:pojecie}
\section{Definicja kategorii}

\begin{definicja}
\emph{Kategorię} $\Cat$ nazywamy parę $(\Cat_0, \Cat_m)$, gdzie:
\begin{enumerate}
\item $\Cat_0$ jest klasą tzw.\ \emph{obiektów}, oznaczanych przez $A, B, C, \ldots$
\item $\Cat_m$ jest klasą tzw.\ \emph{morfizmów} (strzałek), oznaczanych przez $f, g, h, \ldots$
\end{enumerate}
Dla każdej strzałki $f \in \Cat_m$ istnieją $\dom(f) \in \Cat_0$ oraz $\cod(f) \in \Cat_0$, które oznaczamy
\[
f \colon A \longrightarrow B, \qquad \dom(f) = A, \quad \cod(f) = B.
\]
Dla dwóch strzałek $f \colon A \to B$ i $g \colon B \to C$ zdefiniowane jest \emph{złożenie}
\[
g \circ f \colon A \longrightarrow C.
\]
Operacja składania jest \emph{łączna}:
\[
h \circ (g \circ f) = (h \circ g) \circ f
\]
dla $f \colon A \to B$, $g \colon B \to C$, $h \colon C \to D$.

Dla każdego obiektu $A$ istnieje \emph{identyczność} $\id_A \colon A \to A$ taka, że
\[
\id_B \circ f \;=\; f \circ \id_A \;=\; f
\]
dla każdego $f \colon A \to B$.
\end{definicja}

\section{Przykłady kategorii}

\begin{przyklad}[$\Set$]
Kategoria, w której obiektami są wszystkie zbiory, a strzałkami są wszystkie przekształcenia (funkcje), z naturalnym składaniem.
\end{przyklad}

\begin{przyklad}[$\Setf$]
Kategoria, w której obiektami są wszystkie skończone zbiory, a strzałkami są przekształcenia na nich zdefiniowane, z naturalną operacją składania.
\end{przyklad}

\begin{przyklad}[$\Grp$ -- kategoria grup]
Obiektami są wszystkie grupy, morfizmami są homomorfizmy grup, a składanie jest zwykłym składaniem homomorfizmów.
\end{przyklad}

\begin{przyklad}[$\Mon$ -- kategoria monoidów]
Obiektami są wszystkie monoidy, morfizmami są homomorfizmy monoidów, a składanie jest zwykłym składaniem homomorfizmów.
\end{przyklad}

\begin{przyklad}[$\mathbf{Graphs}$]
Kategoria grafów (z homomorfizmami grafów jako strzałkami).
\end{przyklad}

\begin{przyklad}[$\Top$]
Kategoria przestrzeni topologicznych, w której strzałkami są przekształcenia ciągłe.
\end{przyklad}

\begin{przyklad}[$\Pos$]
Obiektami są zbiory częściowo uporządkowane, a strzałkami są przekształcenia zachowujące porządek, z naturalnym składaniem.

Niech $(P, \le)$, $(Q, \le)$ będą obiektami. Strzałka
\[
f \colon (P, \le) \longrightarrow (Q, \le)
\]
to przekształcenie $f \colon P \to Q$ takie, że jeśli $p_1 \le p_2$ w $(P, \le)$, to $f(p_1) \le f(p_2)$ w $(Q, \le)$.

Zwróćmy uwagę na to, że jeśli
\[
f \colon (P, \le) \longrightarrow (Q, \le), \qquad g \colon (Q, \le) \longrightarrow (R, \le)
\]
są przekształceniami zachowującymi porządek, to $g \circ f \colon P \to R$ jest przekształceniem zachowującym porządek. Innymi słowy,
\[
g \circ f \colon (P, \le) \longrightarrow (R, \le)
\]
jest dobrze zdefiniowaną strzałką.
\end{przyklad}

\begin{przyklad}[$\Rel$]
Obiektami są wszystkie zbiory. Dla dwóch obiektów (zbiorów) $A, B$ strzałka $f \colon A \to B$ w $\Rel$ jest relacją $f \subseteq A \times B$. Dla każdego obiektu $A$ mamy
\[
\id_A = \{(a, a) : a \in A\} \subseteq A \times A.
\]
Ponadto, dla strzałek $f \colon A \to B$, $g \colon B \to C$, tzn.\ $f \subseteq A \times B$, $g \subseteq B \times C$, definiujemy złożenie
\[
g \circ f = \bigl\{(a, c) \in A \times C : \exists\, b \in B\ \ (a,b) \in f \ \text{i}\ (b,c) \in g \bigr\}.
\]
Na przykład dla $A = \{a_1, a_2\}$, $B = \{x, y\}$, $C = \{z, t\}$ oraz
\[
f \colon A \to B, \quad f = \{(a_1, y), (a_1, x)\}, \qquad g \colon B \to C, \quad g = \{(x, z), (y, z)\},
\]
mamy
\[
g \circ f = \{(a_1, z)\}.
\]
\end{przyklad}

\begin{przyklad}[Monoid jako kategoria]
Niech $\mathbf{M} = (M, \cdot, 1)$ będzie monoidem. Definiujemy kategorię $\underline{\mathbf{M}}$, w której klasa obiektów składa się z jednego obiektu $*$:
\[
\underline{\mathbf{M}}_0 = \{*\},
\]
a morfizmami są elementy $M$, tzn.\ każdy element $m \in M$ utożsamiamy ze strzałką $m \colon * \to *$. Składaniem w kategorii $\underline{\mathbf{M}}$ jest mnożenie w monoidzie $(M, \cdot, 1)$:
\[
m \circ n \;\stackrel{df}{=}\; m \cdot n,
\]
a identycznością na obiekcie $*$ jest strzałka $1 \colon * \to *$.
\[
\begin{tikzcd}
* \ar[r, "m"] & * \ar[r, "n"] & *
\end{tikzcd}
\]
\end{przyklad}

\begin{przyklad}[Poset jako kategoria]
Niech $\mathbf{P} = (P, \le)$ będzie posetem. Definiujemy kategorię $\underline{\mathbf{P}}$ w następujący sposób:
\begin{enumerate}
\item obiektami $\underline{\mathbf{P}}$ są elementy $P$;
\item dla dwóch elementów $a, b \in P$ istnieje jedna strzałka $a \to b$ wtedy i tylko wtedy, gdy $a \le b$.
\end{enumerate}
Na przykład poset typu ,,diament''
\[
\begin{tikzcd}[row sep=small,column sep=small]
 & 1 \arrow[dl] \arrow[dr] & \\
a \arrow[dr] & & b \arrow[dl] \\
 & 0 &
\end{tikzcd}
\]
wyznacza kategorię o tych samych obiektach i strzałkach $0 \to a,\ 0\to b,\ 0\to 1,\ a\to 1,\ b\to 1$, wzbogaconą o pętle identycznościowe na każdym obiekcie.
\end{przyklad}

\begin{przyklad}[$\Par$ -- przekształcenia częściowe]
Obiekty to wszystkie zbiory. Dla obiektów $A, B$ strzałka $f \colon A \to B$ w $\Par$ to przekształcenie
\[
f \colon A \longrightarrow B + \{\bot\},
\]
gdzie $\bot$ jest \emph{wartością nieokreśloną} (różną od wszystkich elementów $B$), a $+$ oznacza rozłączną sumę zbiorów:
\[
X + Y \;\stackrel{df}{=}\; \{1\} \times X \;\cup\; \{2\} \times Y.
\]
\emph{Interpretacja:} jeśli $f \colon A \to B$ oraz $f(a) = \bot$, to interpretujemy ten wynik jako ,,wartość nieokreśloną''.

Przykładem strzałki w $\Par$ jest dzielenie
\[
\mathrm{div} \colon \mathbb{R} \times \mathbb{R} \longrightarrow \mathbb{R}, \qquad
\bigl(\mathrm{div} \colon \mathbb{R} \times \mathbb{R} \longrightarrow \mathbb{R} + \{\bot\}\bigr),
\]
\[
\mathrm{div}(a, b) = \begin{cases} \dfrac{a}{b}, & b \neq 0, \\ \bot, & \text{w p.p.} \end{cases}
\]

\emph{Składanie w $\Par$.} Dla $f \colon A \to B$, $g \colon B \to C$ złożenie $g \circ f \colon A \to C$ jest dane wzorem
\[
(g \circ f)(a) \;\stackrel{df}{=}\;
\begin{cases}
\bot, & f(a) = \bot, \\
g\bigl(f(a)\bigr), & f(a) \neq \bot.
\end{cases}
\]
\end{przyklad}

\begin{przyklad}[$\Err$ -- kategoria błędów]
Niech $M$ będzie ustalonym zbiorem wiadomości o błędach. Obiekty $\Err$ to wszystkie zbiory. Strzałka $f \colon A \to B$ to przekształcenie
\[
f \colon A \longrightarrow B + M.
\]
\end{przyklad}

\begin{przyklad}[Kategoria $\Trace$ nad monoidem]
Niech $\mathbf{M} = (M, \cdot, 1)$ będzie monoidem. Definiujemy kategorię $\Trace$, w której obiektami są zbiory $A, B, \ldots$, a dla zbiorów $A, B$ strzałka $f \colon A \to B$ jest przekształceniem
\[
f \colon A \longrightarrow M \times B.
\]
Składanie zdefiniowane jest następująco: dla $f \colon A \to B$ (tzn.\ $f \colon A \to M \times B$) oraz $g \colon B \to C$ (tzn.\ $g \colon B \to M \times C$) złożenie
\[
g \circ f \colon A \longrightarrow C \qquad (g \circ f \colon A \to M \times C)
\]
dane jest wzorem
\[
(g \circ f)(a) \;\stackrel{df}{=}\; (m_f \cdot m_g, \, c), \qquad \text{gdzie } f(a) = (m_f, b), \ g(b) = (m_g, c).
\]
Dla każdego zbioru $A$ identyczność
\[
\id_A \colon A \longrightarrow A \qquad (\id_A \colon A \to M \times A)
\]
zadana jest wzorem
\[
\id_A(a) = (1, a).
\]
\end{przyklad}

\section{Funktory}

\begin{definicja}
Niech $\Cat$, $\D$ będą kategoriami. \emph{Funktorem} $F \colon \Cat \to \D$ nazywamy parę przyporządkowań $F = (F_0, F_m)$, tzn.
\[
F_0 \colon \Cat_0 \longrightarrow \D_0, \qquad F_m \colon \Cat_m \longrightarrow \D_m
\]
spełniającą warunki:
\begin{enumerate}
\item dla $f \colon A \to B$ mamy $F_m(f) \colon F_0(A) \to F_0(B)$;
\item $F_m(g \circ f) = F_m(g) \circ F_m(f)$;
\item $F_m(\id_A) = \id_{F_0(A)}$.
\end{enumerate}
\end{definicja}

\begin{uwaga}[Notacja]
Często będziemy pomijali indeksy dolne pisząc $F(A)$ zamiast $F_0(A)$ oraz $F(f)$ zamiast $F_m(f)$. Wtedy warunki (1), (2), (3) brzmią:
\begin{enumerate}
\item $F(f \colon A \to B) = F(f) \colon F(A) \to F(B)$;
\item $F(g \circ f) = F(g) \circ F(f)$ \quad (składanie w $\D$);
\item $F(\id_A) = \id_{F(A)}$ \quad (składanie w $\Cat$).
\end{enumerate}
\end{uwaga}

\subsection{Przykłady funktorów}

\begin{przyklad}[Funktor identycznościowy]
$\Id \colon \Set \to \Set$ lub ogólniej $\Id \colon \Cat \to \Cat$ jest funktorem identycznościowym, tzn.
\[
\Id(A) = A, \qquad \Id(f) = f.
\]
\end{przyklad}

\begin{przyklad}[Funktor $\mathrm{Maybe}$]
$\mathrm{Maybe} \colon \Set \to \Set$. Dla zbioru $B$:
\[
\mathrm{Maybe}\, B \;\stackrel{df}{=}\; B + \{\bot\}.
\]
Dla strzałki $f \colon A \to B$:
\[
\mathrm{Maybe}\, f \colon \mathrm{Maybe}\, A \longrightarrow \mathrm{Maybe}\, B, \qquad
A + \{\bot\} \longrightarrow B + \{\bot\},
\]
\[
(\mathrm{Maybe}\, f)(a) = \begin{cases} f(a), & a \in A, \\ \bot, & a = \bot. \end{cases}
\]
\end{przyklad}

\begin{przyklad}[Funktor $(-)^{*}$ -- wolny monoid]
$(-)^{*} \colon \Set \to \Set$. Dla zbioru $A$ definiujemy
\[
A^{*} \;=\; \text{zbiór wszystkich skończonych ciągów (słów) nad zbiorem } A.
\]
Dla przekształcenia $f \colon A \to B$ definiujemy
\[
f^{*} \colon A^{*} \longrightarrow B^{*}
\]
wzorem
\[
f^{*}(\varepsilon) = \varepsilon, \qquad f^{*}(a_1 \ldots a_n) = f(a_1) \ldots f(a_n).
\]
\end{przyklad}


\chapter{Kategoria produktowa, opposite i slice. Mono- i epimorfizmy}
\section{Podstawowe konstrukcje kategoryjne}

\subsection{Kategoria produktowa}

\begin{definicja}
Niech $\Cat$ i $\D$ będą dwiema kategoriami. \emph{Kategorią produktową} (produktem) $\Cat \times \D$ nazywamy kategorię, w której:
\begin{enumerate}
    \item[(a)] obiektami są pary $(C, D)$, gdzie $C$ jest obiektem $\Cat$ oraz $D$ jest obiektem $\D$;
    \item[(b)] strzałkami z $(C, D)$ do $(C', D')$ w $\Cat \times \D$ są pary
    \[
        (f, g) \colon (C, D) \longrightarrow (C', D'),
    \]
    gdzie $f \colon C \to C'$ w $\Cat$ oraz $g \colon D \to D'$ w $\D$;
    \item[(c)] dla strzałek
    \[
        (f, g) \colon (C, D) \to (C', D'), \qquad (f', g') \colon (C', D') \to (C'', D'')
    \]
    definiujemy złożenie wzorem
    \[
        (f', g') \circ (f, g) \;\stackrel{\mathrm{df}}{=}\; (f' \circ f,\; g' \circ g),
    \]
    gdzie pierwsza składowa to złożenie w $\Cat$, a druga w $\D$. Dodatkowo
    \[
        \id_{(C, D)} \;\stackrel{\mathrm{df}}{=}\; (\id_C, \id_D).
    \]
\end{enumerate}
\end{definicja}

\begin{uwaga}
Należy sprawdzić, że $\Cat \times \D$ jest poprawnie zdefiniowaną kategorią (spełnia aksjomaty łączności i jedynek).
\end{uwaga}

\begin{uwaga}
Dla $\Cat \times \D$ mamy dwa naturalne zadane funktory --- \emph{projekcje}:
\[
\begin{tikzcd}
& \Cat \times \D \arrow[dl, "\pi_1"'] \arrow[dr, "\pi_2"] & \\
\Cat & & \D
\end{tikzcd}
\]
zadane wzorami
\[
    \pi_1(C, D) \stackrel{\mathrm{df}}{=} C, \qquad \pi_1(f, g) \stackrel{\mathrm{df}}{=} f,
\]
oraz analogicznie $\pi_2(C, D) = D$, $\pi_2(f, g) = g$. Należy sprawdzić, że $\pi_1, \pi_2$ są funktorami.
\end{uwaga}

\subsection{Kategoria przeciwna (opposite)}

\begin{definicja}
Niech $\Cat$ będzie kategorią. Definiujemy $\Cat^{\op}$ w następujący sposób:
\begin{enumerate}
    \item[(a)] obiektami $\Cat^{\op}$ są dokładnie obiekty $\Cat$;
    \item[(b)] strzałka $f \colon C \to D$ w $\Cat^{\op}$ jest dokładnie strzałką $f \colon D \to C$ w $\Cat$. Inaczej:
    \[
        \Cat^{\op}(C, D) \stackrel{\mathrm{df}}{=} \Cat(D, C);
    \]
    \item[(c)] dla dwóch strzałek $f \colon C \to D$ i $g \colon D \to E$ w $\Cat^{\op}$ (czyli $f \colon D \to C$ i $g \colon E \to D$ w $\Cat$) definiujemy
    \[
        g \circ_{\Cat^{\op}} f \;\stackrel{\mathrm{df}}{=}\; f \circ_{\Cat} g,
    \]
    czyli złożenie w $\Cat^{\op}$ to złożenie w odwrotnej kolejności w $\Cat$;
    \item[] $\id_A$ w $\Cat^{\op}$ jest $\id_A$ w $\Cat$.
\end{enumerate}
\end{definicja}

\begin{uwaga}
Zauważmy, że $(\Cat^{\op})^{\op} = \Cat$: dwukrotne odwrócenie strzałek przywraca wyjściową kategorię.
\end{uwaga}

\begin{uwaga}
Dla dowolnej kategorii $\Cat$ mamy zadane przyporządkowanie
\[
    \overline{(-)} \colon \Cat \longrightarrow \Cat^{\op},
\]
takie że
\begin{enumerate}
    \item[(a)] $\overline{(A)} \stackrel{\mathrm{df}}{=} A$;
    \item[(b)] $\overline{(f \colon C \to D)} \stackrel{\mathrm{df}}{=} f \colon D \to C$.
\end{enumerate}
To przyporządkowanie nie jest funktorem (zmienia kierunek strzałek).
\end{uwaga}

\subsection{Kategoria slice}

\begin{definicja}
Niech $\Cat$ będzie kategorią i niech $C$ będzie obiektem $\Cat$. \emph{Kategorią slice} (kategorią nad $C$) $\Cat / C$ nazywamy kategorię, w której:
\begin{enumerate}
    \item[(a)] obiektami są strzałki $f \in \Cat$ takie, że $\cod(f) = C$;
    \item[(b)] strzałką z obiektu $f \colon X \to C$ do obiektu $f' \colon X' \to C$ jest morfizm $g \colon X \to X'$ w $\Cat$ taki, że poniższy diagram jest przemienny:
    \[
    \begin{tikzcd}
        X \arrow[rr, "g"] \arrow[dr, "f"'] & & X' \arrow[dl, "f'"] \\
        & C &
    \end{tikzcd}
    \]
    \item[(c)] mając dane dwie strzałki w $\Cat / C$:
    \[
    \begin{tikzcd}
        X \arrow[rr, "g"] \arrow[ddrr, "f"'] & & X' \arrow[rr, "g'"] \arrow[dd, "f'" near start] & & X'' \arrow[ddll, "f''"] \\
        & & & & \\
        & & C & &
    \end{tikzcd}
    \]
    ich złożenie definiujemy wzorem $g' \circ g$. Wówczas
    \[
        f'' \circ g' \circ g = f' \circ g = f,
    \]
    co oznacza, że $g' \circ g$ jest istotnie morfizmem w $\Cat / C$. Identycznością na $f \colon X \to C$ jest $\id_X$.
\end{enumerate}
\end{definicja}

\begin{przyklad}
Niech $\Cat$ będzie kategorią otrzymaną z poseta $(P, \leq)$ oraz niech $p \in P$. Czym jest $P/p$?
\begin{itemize}
    \item Obiektami są strzałki $(x, p)$, gdzie $x \leq p$. Takie obiekty możemy utożsamić z elementami $x$, mianowicie stowarzyszonymi z elementem
    \[
        \{x \colon x \leq p\}.
    \]
    \item Strzałkami w tej kategorii są po prostu nierówności $x \leq x'$ (dla elementów spełniających $x, x' \leq p$).
\end{itemize}
Podsumowując,
\[
    P/p \;=\; \bigl( \mathord{\downarrow}(p),\, \leq \bigr),
\]
gdzie $\mathord{\downarrow}(p) \stackrel{\mathrm{df}}{=} \{x \colon x \leq p\}$ jest \emph{downsetem} elementu $p$.
\end{przyklad}

\begin{uwaga}
Zamieniając w definicji kategorii slice dziedzinę z przeciwdziedziną otrzymujemy definicję \emph{kategorii coslice}, oznaczanej $C / \Cat$. Należy spisać formalną definicję kategorii coslice.
\end{uwaga}

\subsection{Kategorie duże, małe i lokalnie małe}

\begin{uwaga}
Podsumowanie: kategorie duże, małe i lokalnie małe.
\begin{itemize}
    \item Dla ustalonych obiektów $C, D \in \Cat$ kolekcja
    \[
        \{f \colon C \to D \in \Cat\}
    \]
    jest zbiorem $\Longrightarrow$ kategoria jest \emph{lokalnie mała}.
    \item Jeżeli kolekcja obiektów i kolekcja strzałek są zbiorami, kategoria jest \emph{mała}; w przeciwnym razie jest \emph{duża}. Kategoria duża może być lokalnie mała --- np.\ $\Set$ jest duża (klasa wszystkich zbiorów nie jest zbiorem), ale lokalnie mała.
\end{itemize}
\end{uwaga}

\section{Struktury abstrakcyjne: izo, mono, epi}

\begin{definicja}
Strzałkę $f \colon A \to B$ w $\Cat$ nazywamy \emph{izomorfizmem} (izo), jeśli istnieje strzałka $g \colon B \to A$ w $\Cat$ taka, że
\[
    g \circ f = \id_A \quad \text{oraz} \quad f \circ g = \id_B.
\]
\end{definicja}

\begin{przyklad}
Izomorfizmami w $\Set$ są dokładnie bijekcje. W $\Grp$ izomorfizmy są standardowymi izomorfizmami grup; analogicznie w innych kategoriach struktur algebraicznych.
\end{przyklad}

\begin{twierdzenie}
Jeżeli $f \colon A \to B$ i $f' \colon B \to C$ są izomorfizmami w $\Cat$, to $f' \circ f$ jest izomorfizmem.
\end{twierdzenie}

\begin{definicja}
W kategorii $\Cat$ strzałka $f \colon A \to B$ nazywana jest
\begin{enumerate}
    \item[(a)] \emph{monomorfizmem} (mono), jeśli dla każdej pary $g, h \colon C \to A$ mamy
    \[
        f \circ g = f \circ h \;\Longrightarrow\; g = h;
    \]
    \item[(b)] \emph{epimorfizmem} (epi), jeśli dla każdej pary $i, j \colon B \to D$ mamy
    \[
        i \circ f = j \circ f \;\Longrightarrow\; i = j.
    \]
\end{enumerate}
\end{definicja}

\begin{uwaga}
Oznaczenia:
\[
    A \hookrightarrow B \;\text{---}\; \text{monomorfizm}, \qquad A \twoheadrightarrow B \;\text{---}\; \text{epimorfizm}.
\]
\end{uwaga}

\begin{stwierdzenie}
Funkcja $f \colon A \to B$ w $\Set$ jest mono wtedy i tylko wtedy, gdy jest różnowartościowa.
\end{stwierdzenie}

\begin{uwaga}
Zwróćmy uwagę na to, że
\[
    \text{strzałka } f \colon A \to B \text{ jest epi w } \Cat
    \;\Longleftrightarrow\;
    \text{strzałka } f \colon B \to A \text{ jest mono w } \Cat^{\op}.
\]
\end{uwaga}


\chapter{Obiekty początkowe i końcowe. Produkty}
\section{Monomorfizmy, epimorfizmy, izomorfizmy}

Przypomnijmy najpierw pojęcia izomorfizmu, mono- i epimorfizmu wprowadzone pod koniec poprzedniego rozdziału; zobaczymy je teraz w działaniu na przykładach.

\begin{definicja}
Niech $f\colon A\to B$ będzie strzałką w kategorii $\C$. Strzałkę $f$ nazywamy \emph{izomorfizmem}, jeśli istnieje $g\colon B\to A$ w $\C$ taka, że
\[
g\circ f = \id_A \quad\text{oraz}\quad f\circ g = \id_B.
\]
\end{definicja}

\begin{definicja}
Strzałkę $f\colon A\to B$ nazywamy \emph{monomorfizmem} (\emph{mono}), jeśli dla każdych $g,h\colon C\to A$ w $\C$
\[
f\circ g = f\circ h \;\Longrightarrow\; g=h,
\]
tzn.\ diagram
\[
\begin{tikzcd}
C \arrow[r, shift left, "g"] \arrow[r, shift right, "h"'] & A \arrow[r, "f"] & B
\end{tikzcd}
\]
daje implikację $f\circ g = f\circ h \Rightarrow g=h$.
\end{definicja}

\begin{definicja}
Strzałkę $f\colon A\to B$ nazywamy \emph{epimorfizmem} (\emph{epi}), jeśli dla każdych $i,j\colon B\to C$ w $\C$
\[
i\circ f = j\circ f \;\Longrightarrow\; i=j.
\]
\end{definicja}

\begin{stwierdzenie}
Funkcja $f\colon A\to B$ jest mono w $\Set$ wtedy i tylko wtedy, gdy jest różnowartościowa.
\end{stwierdzenie}

\begin{stwierdzenie}
$f\colon A\to B$ jest mono w $\C$ wtedy i tylko wtedy, gdy $f\colon B\to A$ jest epi w $\C^{\op}$.
\end{stwierdzenie}

\begin{proof}
Niech $f\colon A\to B$ będzie mono w $\C$. W szczególności oznacza to, że $f\colon B\to A$ jest strzałką w $\C^{\op}$. Weźmy $i,j\colon A\to C$ w $\C^{\op}$ i załóżmy, że w $\C^{\op}$ zachodzi
\[
i\circ_{\C^{\op}} f = j\circ_{\C^{\op}} f,
\]
gdzie $f\colon A\to B$, $i,j\colon C\to A$ w $\C$. Z definicji składania w $\C^{\op}$ wynika, że w $\C$ zachodzi
\[
f\circ i = f\circ j,
\]
gdzie $f\colon A\to B$ oraz $i,j\colon C\to A$ w $\C$. Ponieważ $f$ jest mono w $\C$, dostajemy $i=j$ w $\C$. Ale to znaczy, że $i=j$ w $\C^{\op}$. Stąd $f$ jest epi w $\C^{\op}$. Implikacja w przeciwną stronę jest analogiczna.
\end{proof}

\subsection{Przykłady}

\begin{przyklad}
W kategorii $\Set$ mono $=$ ,,różnowartościowe'', a epi $=$ ,,na''.
\end{przyklad}

\begin{przyklad}
Niech $\C$ będzie kategorią indukowaną przez poset $(P,\leq)$. Weźmy $p,q\in P$. Wtedy dla strzałek
\[
\begin{tikzcd}
g,h\colon r \arrow[r] & p \arrow[r, "f"] & q
\end{tikzcd}
\]
warunek $f\circ g = f\circ h$ daje $g=h$ (gdyż wszystkie strzałki między tą samą parą obiektów są równe).
\end{przyklad}

\begin{wniosek}
Każda strzałka $f\colon p\leq q$ jest mono (i epi) w kategorii indukowanej przez poset $(P,\leq)$.
\end{wniosek}

\begin{przyklad}
Rozważmy kategorię $\Mon$ -- kategorię monoidów z homomorfizmami monoidów jako strzałkami. W tej kategorii rozważmy strzałkę
\[
f\colon (\mathbb{N},+,0) \longrightarrow (\mathbb{Z},+,0), \qquad f(i) = i.
\]
Ta strzałka jest:
\begin{enumerate}
\item mono,
\item epimorfizmem.
\end{enumerate}

\emph{Dowód (1).} Strzałka $f$ jest różnowartościowa, więc jest mono w $\Set$ -- a zatem także w $\Mon$.

\emph{Dowód (2).} Musimy pokazać, że dla każdej pary
\[
i,j\colon (\mathbb{Z},+,0) \longrightarrow M
\]
(gdzie $M=(M,\cdot,1)$ jest monoidem) zachodzi
\[
i\circ f = j\circ f \;\Longrightarrow\; i=j.
\]
Załóżmy, że $i\circ f = j\circ f$. Z tej równości wynika w szczególności, że
\[
i|_{\mathbb{N}} = j|_{\mathbb{N}}.
\]
Wystarczy więc pokazać, że $i(-n) = j(-n)$ dla każdego $n\in\mathbb{N}$. Z faktu, że $i$ jest homomorfizmem mamy
\[
1 = i(0) = i(n + (-n)) = i(n)\cdot i(-n)
\qquad\text{oraz}\qquad
1 = i(0) = i((-n) + n) = i(-n)\cdot i(n),
\]
czyli $i(-n)$ jest obustronną odwrotnością elementu $i(n)$. W monoidzie element ma co najwyżej jedną obustronną odwrotność: jeśli $x\cdot m = 1 = m\cdot y$, to $x = x\cdot(m\cdot y) = (x\cdot m)\cdot y = y$. Analogicznie $j(-n)$ jest obustronną odwrotnością $j(n)$. Z $i|_{\mathbb{N}}=j|_{\mathbb{N}}$ wynika $i(n)=j(n)$, więc z jednoznaczności odwrotności $i(-n) = j(-n)$, a zatem $i=j$. \qed
\end{przyklad}

\begin{uwaga}
Powyższy przykład pokazuje, że strzałka będąca jednocześnie mono i epi nie musi być izomorfizmem ($f$ nie jest ,,na'', więc nie ma odwrotności). W $\Set$ tak być nie może: mono $+$ epi $=$ różnowartościowa i ,,na'' $=$ bijekcja $=$ izo. Inny przykład: w $\Top$ włożenie $\mathbb{Q} \hookrightarrow \mathbb{R}$ jest mono (jest różnowartościowe) i epi (dwa przekształcenia ciągłe o wartościach w przestrzeni Hausdorffa, zgodne na gęstym podzbiorze, są równe), ale nie jest homeomorfizmem.
\end{uwaga}

\section{Obiekty początkowe i końcowe}

W $\Set$:
\[
\begin{tikzcd}
\emptyset \arrow[r, "i"] & A \arrow[r] & \mathbf{1}
\end{tikzcd}
\]
$\emptyset$ jest obiektem początkowym, $\mathbf{1}$ (zbiór jednoelementowy) jest obiektem końcowym.

\begin{definicja}
Niech $\C$ będzie kategorią. Obiekt $0\in\C$ nazywamy \emph{początkowym}, jeśli dla każdego obiektu $A\in\C$ istnieje dokładnie jedna strzałka
\[
i_A\colon 0 \longrightarrow A.
\]
\end{definicja}

\begin{definicja}
Obiekt $\mathbf{1}\in\C$ nazywamy \emph{końcowym}, jeśli dla każdego $A\in\C$ istnieje dokładnie jedna strzałka
\[
!_A\colon A \longrightarrow \mathbf{1}.
\]
\end{definicja}

\begin{uwaga}
$0\in\C$ jest obiektem początkowym w $\C$ wtedy i tylko wtedy, gdy $0$ jest końcowym w $\C^{\op}$.
\end{uwaga}

\begin{stwierdzenie}
Obiekt początkowy (końcowy) jest wyznaczony jednoznacznie z dokładnością do izomorfizmu.
\end{stwierdzenie}

\begin{proof}
Załóżmy, że $0$, $0'$ są początkowe w kategorii $\C$. Wtedy istnieją jednoznacznie wyznaczone strzałki
\[
\begin{tikzcd}
0 \arrow[r, shift left, "u"] & 0' \arrow[l, shift left, "v"]
\end{tikzcd}
\]
spełniające:
\begin{enumerate}
\item $v\circ u = \id_0$ (gdyż obie strzałki $0\to 0$ muszą być równe, a $\id_0$ jest jedyną),
\item $u\circ v = \id_{0'}$ (analogicznie).
\end{enumerate}
Stąd $0\cong 0'$. Dowód dla obiektu końcowego jest analogiczny.
\end{proof}

\subsection{Przykłady}

\begin{przyklad}
W $\Set$: obiekt początkowy $=\emptyset$, obiekt końcowy $=$ dowolny zbiór jednoelementowy.
\end{przyklad}

\begin{przyklad}
W $\Grp$ grupa trywialna $\{e\}$ jest jednocześnie obiektem początkowym (jedyny homomorfizm $\{e\} \to G$ posyła $e$ na element neutralny) i końcowym (jedyny homomorfizm $G \to \{e\}$ jest stały). Obiekt o obu tych własnościach nazywamy \emph{obiektem zerowym}. W $\Set$ obiekty początkowy i końcowy są natomiast różne: $\emptyset \neq \{*\}$.
\end{przyklad}

\begin{przyklad}
Niech $\C$ będzie kategorią indukowaną przez poset $(P,\leq)$. Czym (jeśli istnieją) są obiekty początkowe i końcowe w $\C$?

\emph{Obiekt początkowy:} $0\leq p$ dla każdego $p\in P$. Widzimy, że $0$ jest elementem najmniejszym w $(P,\leq)$. Np.\ w
\[
\begin{tikzcd}
\bullet & & \bullet
\end{tikzcd}
\]
nie ma obiektu początkowego.

\emph{Obiekt końcowy:} $p \leq \mathbf{1}$ dla każdego $p\in P$, czyli $\mathbf{1}$ jest elementem największym w $(P,\leq)$.
\end{przyklad}

\begin{przyklad}
Rozważmy kategorię algebr typu $(1,0)$. Obiektami w tej kategorii są
\[
(A,s,a), \qquad \text{gdzie } s\colon A\to A,\; a\in A.
\]
Strzałkami w tej kategorii są
\[
f\colon (A,s,a) \longrightarrow (B,s',b),
\]
gdzie $f\colon A\to B$ jest funkcją spełniającą:
\begin{enumerate}
\item dla każdego $x\in A$: $s'(f(x)) = f(s(x))$,
\item $f(a)=b$.
\end{enumerate}
\end{przyklad}

\begin{stwierdzenie}
W tej kategorii obiektem początkowym jest algebra
\[
(\mathbb{N},\mathrm{succ},0), \qquad \mathrm{succ}(n)\stackrel{\mathrm{df}}{=} n+1.
\]
\end{stwierdzenie}

\begin{proof}[Szkic dowodu]
Niech $(B, s', b)$ będzie dowolną algebrą typu $(1,0)$. Strzałka $f\colon (\mathbb{N},\mathrm{succ},0)\to (B,s',b)$ musi spełniać
\[
f(0) = b, \qquad f(n+1) = f(\mathrm{succ}(n)) = s'(f(n)),
\]
a te dwa warunki \emph{definiują} $f$ rekurencyjnie -- strzałka więc istnieje, a jej jednoznaczność pokazujemy przez indukcję po $n$. Inicjalność $(\mathbb{N},\mathrm{succ},0)$ to zatem dokładnie zasada definiowania przez rekursję.
\end{proof}

\section{Produkty}

W $\Set$:
\[
A\times B = \{(a,b) : a\in A,\; b\in B\}.
\]
Zauważmy, że każdy produkt kartezjański $A\times B$ jest w parze z dwoma rzutowaniami
\[
\begin{tikzcd}
A & A\times B \arrow[l, "\pi_1"'] \arrow[r, "\pi_2"] & B
\end{tikzcd}
\]
gdzie $\pi_1(a,b) = a$, $\pi_2(a,b)=b$. Ponadto dla dowolnego zbioru $X$ wraz ze strzałkami $c_1\colon X\to A$, $c_2\colon X\to B$ istnieje jednoznacznie wyznaczona strzałka $X\to A\times B$ dana wzorem $x\mapsto (c_1(x),c_2(x))$.

\begin{definicja}
Niech $\C$ będzie kategorią. \emph{Diagramem produktu} dla obiektów $A,B\in\C$ nazywamy obiekt $P$ wraz ze strzałkami
\[
\begin{tikzcd}
A & P \arrow[l, "p_1"'] \arrow[r, "p_2"] & B
\end{tikzcd}
\]
takimi, że dla każdego obiektu $X\in\C$ wraz ze strzałkami
\[
\begin{tikzcd}
A & X \arrow[l, "x_1"'] \arrow[r, "x_2"] & B
\end{tikzcd}
\]
istnieje jednoznaczna strzałka $u\colon X\to P$ czyniąca następujący diagram przemiennym:
\[
\begin{tikzcd}
& X \arrow[ld, "x_1"'] \arrow[rd, "x_2"] \arrow[d, dashed, "\exists! u"] & \\
A & P \arrow[l, "p_1"] \arrow[r, "p_2"'] & B
\end{tikzcd}
\]
\end{definicja}

\begin{definicja}
Jeśli $P$ wraz z $A\xleftarrow{p_1} P\xrightarrow{p_2} B$ jest diagramem produktowym dla obiektów $A,B$, to $P$ nazywamy \emph{produktem} $A$ oraz $B$.
\end{definicja}

\begin{stwierdzenie}
Produkty są określone jednoznacznie z dokładnością do izomorfizmu.
\end{stwierdzenie}

\begin{proof}
Załóżmy, że
\[
\begin{tikzcd}
A & P \arrow[l, "p_1"'] \arrow[r, "p_2"] & B
\end{tikzcd}
\qquad\text{oraz}\qquad
\begin{tikzcd}
A & Q \arrow[l, "q_1"'] \arrow[r, "q_2"] & B
\end{tikzcd}
\]
są diagramami produktowymi dla $A,B$. $Q$ jest produktem, więc istnieje jednoznacznie określona strzałka $i\colon P\to Q$ taka, że
\[
q_1\circ i = p_1 \quad\text{oraz}\quad q_2\circ i = p_2.
\]
Analogicznie, ponieważ $P$ jest produktem, istnieje jednoznacznie określona strzałka $j\colon Q\to P$ taka, że
\[
q_1 = p_1\circ j, \qquad q_2 = p_2\circ j.
\]
Stąd $p_1\circ j\circ i = p_1$ oraz $p_2\circ j\circ i = p_2$. Z jednoznaczności strzałki czyniącej diagram produktowy przemiennym (a $\id_P$ także go uzupełnia) mamy $j\circ i = \id_P$. Analogicznie $i\circ j = \id_Q$.
\end{proof}


\chapter{Skończone (ko)granice: produkty, koprodukty, ekwalizatory i pullbacki}
\section{Produkty, dualność i koprodukty}
Niech $\C$ będzie kategorią.

\begin{definicja}
\emph{Diagramem produktowym} dla $A, B \in \C$ nazywamy obiekt $P$ wraz ze strzałkami
\[
\begin{tikzcd}
A & P \arrow[l, "p_1"'] \arrow[r, "p_2"] & B
\end{tikzcd}
\]
takimi, że dla każdego obiektu $X \in \C$ wraz ze strzałkami
\[
\begin{tikzcd}
A & X \arrow[l, "x_1"'] \arrow[r, "x_2"] & B
\end{tikzcd}
\]
istnieje jednoznaczne $u \colon X \to P$, które czyni diagram
\[
\begin{tikzcd}
& X \arrow[dl, "x_1"'] \arrow[dr, "x_2"] \arrow[d, "u" description, dashed] & \\
A & P \arrow[l, "p_1"] \arrow[r, "p_2"'] & B
\end{tikzcd}
\]
przemiennym. Piszemy $P = A \times B$ oraz $u = \langle x_1, x_2 \rangle$.
\end{definicja}

\begin{przyklad}
\leavevmode
\begin{enumerate}
\item W kategorii $\Set$ dla dwóch zbiorów $A, B \in \Set$ ich diagramem produktowym jest
\[
\begin{tikzcd}
A & A \times B \arrow[l, "\pi_1"'] \arrow[r, "\pi_2"] & B
\end{tikzcd}
\qquad A \times B = \{(a, b) \mid a \in A,\ b \in B\}.
\]

\item W omawianych wcześniej kategoriach algebr ($\Mon$, $\Grp$, $\Ring$ itd.) dla każdych dwóch obiektów $A, B$ istnieje produkt $A \times B$ zdefiniowany w naturalny sposób (np. dla $\Grp$ -- produkt grup).

\item Niech $\mathcal{P}$ będzie kategorią indukowaną przez $(P, \leq)$. Czym jest produkt $p \times q$ dla $p, q \in P$?
\[
\begin{tikzcd}
& p \times q \arrow[dl] \arrow[dr] & \\
p & & q
\end{tikzcd}
\]
\textbf{Obs.}
\begin{enumerate}
\item[a)] $p \times q \leq p$ oraz $p \times q \leq q$;
\item[b)] $\forall x$ t.\,że $x \leq p$, $x \leq q$ mamy $x \leq p \times q$.
\end{enumerate}
Widać, że $p \times q = p \wedge q$.
\end{enumerate}
\end{przyklad}

\begin{uwaga}
W kategorii indukowanej przez antyłańcuch $p \not\leq q$, $q \not\leq p$ nie ma produktu $p \times q$.
\end{uwaga}

\subsection*{Własności kategorii z produktem}

Niech $\C$ będzie kategorią, w której dla każdych dwóch obiektów istnieje ich diagram produktowy. Dodatkowo zakładamy, że mamy następujący diagram
\[
\begin{tikzcd}
A & A \times A' \arrow[l, "p_1"'] \arrow[r, "p_2"] & A' \\
B \arrow[u, "f"] & B \times B' \arrow[l, "q_1"] \arrow[r, "q_2"'] \arrow[u, dashed, "\exists! u"] & B' \arrow[u, "f'"']
\end{tikzcd}
\]
Oznaczamy $u = f \times f'$.

W podobny sposób możemy zdefiniować dla trzech obiektów $A, B, C$ ich produkt
\[
\begin{tikzcd}
& A \times B \times C \arrow[dl] \arrow[d] \arrow[dr] & \\
A & B & C
\end{tikzcd}
\quad (\ast)
\]

\begin{stwierdzenie}
Jeśli $\C$ jest kategorią z produktami binarnymi, to w $\C$ istnieją produkty ternarne $(\ast)$. Dla $A, B, C \in \C$ $(A \times B) \times C$ jest produktem ternarnym $A, B, C$.
\end{stwierdzenie}

\begin{definicja}
Mówimy, że kategoria $\C$ ma \emph{wszystkie skończone produkty}, jeśli
\begin{enumerate}
\item ma obiekt końcowy;
\item dla każdych $A, B \in \C$ istnieje nad nimi diagram produktowy.
\end{enumerate}
\end{definicja}

\subsection*{Więcej o dualności}

Niech $\C$ to:
\begin{itemize}
\item $\C_0$: $A, B, C$ -- obiekty;
\item $\C_1$: $f, g, h$ -- strzałki;
\item operacje: $\dom(f)$, $\cod(f)$, $\id_A$, $g \circ f$,
\end{itemize}
spełniające aksjomaty:
\[
\dom(\id_A) = A, \qquad \cod(\id_A) = A,
\]
\[
f \circ \id_{\dom(f)} = f, \qquad \id_{\cod(f)} \circ f = f,
\]
\[
\dom(g \circ f) = \dom(f), \qquad \cod(g \circ f) = \cod(g),
\]
\[
h \circ (g \circ f) = (h \circ g) \circ f.
\]

Niech $\Sigma$ będzie dowolnym zdaniem wyrażonym w podstawowym języku teorii kategorii. Możemy zdefiniować zdanie $\Sigma^{\op}$ dualne do $\Sigma$ zamieniając:
\begin{enumerate}
\item każde wystąpienie $f \circ g$ na $g \circ f$;
\item $\cod$ na $\dom$;
\item $\dom$ na $\cod$.
\end{enumerate}

Widać, że $\Sigma^{\op}$ też jest poprawnym zdaniem w TK. Zachodzi stąd: \emph{Dla każdego poprawnego zdania $\Sigma$, jeśli $\Sigma$ wynika z aksjomatów TK, to $\Sigma^{\op}$ też.}

\subsection*{Koprodukty}

Niech $\C$ będzie kategorią.

\begin{definicja}
Dla $A, B \in \C$ \emph{diagramem koproduktowym} nad $A, B$ nazywamy obiekt $Q$ wraz ze strzałkami
\[
\begin{tikzcd}
A \arrow[r, "q_1"] & Q & B \arrow[l, "q_2"']
\end{tikzcd}
\]
takimi, że dla każdego obiektu $Z$ oraz $A \xrightarrow{z_1} Z \xleftarrow{z_2} B$ istnieje jednoznaczne $u$ czyniące następujący diagram przemiennym
\[
\begin{tikzcd}
A \arrow[r, "q_1"] \arrow[dr, "z_1"'] & Q \arrow[d, dashed, "\exists! u"] & B \arrow[l, "q_2"'] \arrow[dl, "z_2"] \\
& Z &
\end{tikzcd}
\]
Oznaczenie: $u = [z_1, z_2]$, $Q = A + B$.
\end{definicja}

\begin{przyklad}
\leavevmode
\begin{enumerate}
\item Czym są koprodukty (jeśli istnieją) w $\Set$? Weźmy
\[
A + B = A \times \{1\} \cup B \times \{2\} \qquad (\text{rozłączna suma kopii}),
\]
oraz
\[
i_1 \colon A \to A + B, \quad a \mapsto (a, 1), \qquad i_2 \colon B \to A + B, \quad b \mapsto (b, 2).
\]
Wówczas
\[
\begin{tikzcd}
A \arrow[r, "i_1"] & A + B & B \arrow[l, "i_2"']
\end{tikzcd}
\]
jest diagramem koproduktowym nad $A, B$.

\item W kategorii $\mathcal{P}$ indukowanej przez poset $(P, \leq)$:
\[
\begin{tikzcd}
p \arrow[r] & p + q & q \arrow[l]
\end{tikzcd}
\]
Widać, że $p + q = p \vee q$.
\end{enumerate}
\end{przyklad}

Niech dany jest następujący diagram
\[
\begin{tikzcd}
A \arrow[r, "a_1"] \arrow[d, "f"'] & A + A' \arrow[d, dashed, "\exists! u"] & A' \arrow[l, "a_2"'] \arrow[d, "f'"] \\
B \arrow[r, "b_1"'] & B + B' & B' \arrow[l, "b_2"]
\end{tikzcd}
\]
Dla każdej pary $f \colon A \to B$, $f' \colon A' \to B'$ istnieje jednoznaczna strzałka
\[
f + f' \colon A + A' \to B + B'
\]
(o ile $\C$ jest kategorią z koproduktami binarnymi).

\begin{stwierdzenie}
Koprodukty są określone jednoznacznie z dokładnością do izomorfizmu.
\end{stwierdzenie}
\begin{proof}
Wynika z zasady dualności.
\end{proof}

\begin{definicja}
Mówimy, że $\C$ ma \emph{wszystkie skończone koprodukty}, jeśli:
\begin{enumerate}
\item $\C$ ma obiekt początkowy;
\item $\C$ ma wszystkie binarne koprodukty.
\end{enumerate}
\end{definicja}

\section{Ekwalizatory, koekwalizatory i pullbacki}
\subsection*{Ekwalizatory}

Niech $\C$ będzie kategorią.

\begin{definicja}
Dla pary strzałek $f, g \colon A \to B \in \C$ mówimy, że $e\colon E\to A$ jest \emph{ekwalizatorem} $f$ i $g$, jeśli:
\begin{enumerate}
\item diagram
\[\begin{tikzcd}
E \arrow[r,"e"] & A \arrow[r,shift left,"f"] \arrow[r,shift right,swap,"g"] & B
\end{tikzcd}\]
jest przemienny, tj.\ $f\circ e = g\circ e$,
\item dla dowolnego $z\colon Z\to A$ takiego, że $f\circ z = g\circ z$, istnieje dokładnie jedna strzałka $u\colon Z\to E$ czyniąca następujący diagram przemiennym:
\[\begin{tikzcd}
E \arrow[r,"e"] & A \arrow[r,shift left,"f"] \arrow[r,shift right,swap,"g"] & B \\
Z \arrow[ur,swap,"z"] \arrow[u,dashed,"\exists !\, u"] & &
\end{tikzcd}\]
\end{enumerate}
\end{definicja}

\begin{przyklad}
W kategorii $\Set$, mając dane dwa przekształcenia $f,g\colon A\to B$, ich ekwalizatorem będzie
\[
\{x\in A \mid f(x)=g(x)\} \hookrightarrow A,
\]
gdzie strzałka jest zanurzeniem (inkluzją).
\end{przyklad}

\paragraph{Dowód.}
Dla każdego $z\colon Z\to A$ takiego, że $f\circ z = g\circ z$, mamy: dla $y\in Z$ zachodzi $f(z(y)) = g(z(y))$, czyli element $x:=z(y)\in A$ spełnia $f(x)=g(x)$. Definiujemy
\[
u\colon Z \to \{x\in A\mid f(x)=g(x)\}, \qquad u(y) := z(y).
\]
Wtedy diagram
\[\begin{tikzcd}
\{x\in A \mid f(x)=g(x)\} \arrow[r,hook] & A \arrow[r,shift left,"f"] \arrow[r,shift right,swap,"g"] & B \\
Z \arrow[ur,swap,"z"] \arrow[u,"u"] & &
\end{tikzcd}\]
jest przemienny. Przekształcenie $u$ jest jednoznaczne. \qed

\begin{stwierdzenie}
Jeśli $e\colon E\to A$ jest ekwalizatorem pary $A\rightrightarrows B$, to $e$ jest mono.
\end{stwierdzenie}

\begin{przyklad}
Rozważmy w $\Set$ podzbiór $U\subseteq A$ oraz diagram
\[\begin{tikzcd}
? \arrow[r,"?"] & A \arrow[r,shift left,"\mathrm{const}_1"] \arrow[r,shift right,swap,"\chi_U"] & 2=\{0,1\}
\end{tikzcd}\]
gdzie
\[
\chi_U(a)=\begin{cases} 1 & a\in U,\\ 0 & \text{w p.p.},\end{cases} \qquad \mathrm{const}_1(a)=1.
\]
Czym jest ekwalizator $\mathrm{const}_1$ i $\chi_U$? Warunek $\chi_U(a) = \mathrm{const}_1(a) = 1$ oznacza dokładnie $a\in U$, więc ekwalizatorem jest inkluzja $U \hookrightarrow A$. W szczególności \emph{każdy} podzbiór $U\subseteq A$ jest ekwalizatorem pewnej pary strzałek o wartościach w $2=\{0,1\}$.
\end{przyklad}

\subsection*{Koekwalizatory}

\begin{definicja}
Dla pary strzałek $f, g \colon A\to B \in \C$ \emph{koekwalizatorem} $f$ i $g$ nazywamy strzałkę $q\colon B\to Q$ taką, że:
\begin{enumerate}
\item $q\circ f = q\circ g$, tj.\ diagram
\[\begin{tikzcd}
A \arrow[r,shift left,"f"] \arrow[r,shift right,swap,"g"] & B \arrow[r,"q"] & Q
\end{tikzcd}\]
jest przemienny,
\item dla każdego $z\colon B\to Z$ takiego, że $z\circ f = z\circ g$, istnieje dokładnie jedna strzałka $u\colon Q\to Z$ taka, że poniższy diagram jest przemienny:
\[\begin{tikzcd}
A \arrow[r,shift left,"f"] \arrow[r,shift right,swap,"g"] & B \arrow[r,"q"] \arrow[dr,swap,"z"] & Q \arrow[d,dashed,"\exists !\, u"] \\
& & Z
\end{tikzcd}\]
\end{enumerate}
\end{definicja}

\begin{stwierdzenie}
Jeśli $q\colon B\to Q$ jest koekwalizatorem pewnej pary strzałek, to $q$ jest epi.
\end{stwierdzenie}

\begin{przyklad}
Niech $R\subseteq X\times X$ będzie relacją binarną na $X$. Rozważmy diagram
\[\begin{tikzcd}
R \arrow[r,shift left,"\pi_1"] \arrow[r,shift right,swap,"\pi_2"] & X
\end{tikzcd}\]
Jeśli $R$ jest relacją równoważności, to koekwalizatorem tej pary rzutowań jest
\[\begin{tikzcd}
R \arrow[r,shift left,"\pi_1"] \arrow[r,shift right,swap,"\pi_2"] & X \arrow[r,"\mathrm{nat}\,R"] & X/R,
\end{tikzcd}\]
gdzie
\[
X/R = \{\,[x]_R \mid x\in X\,\}
\]
(klasy abstrakcji relacji $R$),
\[
\mathrm{nat}\,R\colon X \twoheadrightarrow X/R, \qquad \mathrm{nat}\,R(x) \stackrel{\mathrm{df}}{=} [x]_R.
\]
\end{przyklad}

\subsection*{Pullback}

\begin{definicja}
\emph{Pullbackiem} (cofnięciem) pary strzałek
\[\begin{tikzcd}
& B \arrow[d,"g"] \\
A \arrow[r,"f"] & C
\end{tikzcd}\]
nazywamy obiekt $P$ wraz ze strzałkami $p_1\colon P\to A$, $p_2\colon P\to B$ takimi, że
\begin{enumerate}
\item diagram
\[\begin{tikzcd}
P \arrow[r,"p_2"] \arrow[d,swap,"p_1"] & B \arrow[d,"g"] \\
A \arrow[r,"f"] & C
\end{tikzcd}\]
jest przemienny,
\item dla każdego obiektu $Z$ wraz ze strzałkami $z_1\colon Z\to A$ oraz $z_2\colon Z\to B$ takimi, że $f\circ z_1 = g\circ z_2$, istnieje dokładnie jedna strzałka $u\colon Z\to P$ czyniąca następujący diagram przemiennym:
\[\begin{tikzcd}
Z \arrow[drr,bend left,"z_2"] \arrow[ddr,bend right,swap,"z_1"] \arrow[dr,dashed,"\exists !\, u"] & & \\
& P \arrow[r,"p_2"] \arrow[d,swap,"p_1"] & B \arrow[d,"g"] \\
& A \arrow[r,swap,"f"] & C
\end{tikzcd}\]
\end{enumerate}
\end{definicja}

\begin{uwaga}
Często używać będziemy notacji ,,produktowej'' w kontekście pullbacków. Dokładniej, obiekt pullback strzałek $A\xrightarrow{f} C \xleftarrow{g} B$ oznaczać będziemy przez $A\times_C B$. Morfizmy pullbacka często będziemy też oznaczać przez $\pi_1,\pi_2$:
\[\begin{tikzcd}
A\times_C B \arrow[r,"\pi_2"] \arrow[d,swap,"\pi_1"] & B \arrow[d,"g"] \\
A \arrow[r,swap,"f"] & C
\end{tikzcd}\]
\end{uwaga}

\begin{przyklad}
W $\Set$ pullbackiem pary $A\xrightarrow{f} C \xleftarrow{g} B$ jest
\[
A\times_C B \;=\; \{(a,b)\in A\times B \;:\; f(a) = g(b)\}
\]
wraz z obcięciami rzutowań. Dwa ważne przypadki szczególne:
\begin{enumerate}
\item \emph{Przeciwobraz.} Dla inkluzji $i_U \colon U \hookrightarrow C$ podzbioru $U\subseteq C$ mamy
\[
A \times_C U \;=\; \{(a,u) : f(a) = u\} \;\cong\; f^{-1}(U),
\]
czyli pullback wzdłuż $f$ ,,cofa'' podzbiór $U$ do jego przeciwobrazu (por.\ zadanie~\ref{zad:pullback-przeciwobraz}). Stąd właśnie nazwa \emph{cofnięcie}.
\item \emph{Przecięcie.} Dla inkluzji $U \hookrightarrow X \hookleftarrow V$ dwóch podzbiorów $U, V \subseteq X$ pullbackiem jest $U \cap V$ wraz z inkluzjami.
\end{enumerate}
\end{przyklad}

\begin{stwierdzenie}
Niech $\C$ będzie kategorią z binarnymi produktami i ekwalizatorami. Dla danej pary $A\xrightarrow{f} C \xleftarrow{g} B$:
\begin{enumerate}
\item weźmy produkt binarny $A\times B$ wraz z $\pi_1\colon A\times B\to A$, $\pi_2\colon A\times B\to B$,
\item rozważmy ekwalizator $E \xhookrightarrow{e} A\times B$ pary
\[\begin{tikzcd}
A\times B \arrow[r,shift left,"f\circ \pi_1"] \arrow[r,shift right,swap,"g\circ \pi_2"] & C.
\end{tikzcd}\]
\end{enumerate}
Wtedy obiekt $E$ wraz z $\pi_1\circ e\colon E\to A$ oraz $\pi_2\circ e\colon E\to B$ jest pullbackiem $A\xrightarrow{f} C \xleftarrow{g} B$.
\end{stwierdzenie}

\paragraph{Dowód.}
Widzimy, że $E$ wraz z $\pi_1\circ e$ oraz $\pi_2\circ e$ czyni kwadrat $f\circ(\pi_1\circ e) = g\circ(\pi_2\circ e)$ przemiennym -- wynika to z faktu, że $e$ jest ekwalizatorem $f\circ\pi_1$ oraz $g\circ\pi_2$.

Weźmy obiekt $Z$ wraz z $z_1\colon Z\to A$, $z_2\colon Z\to B$ takie, że $f\circ z_1 = g\circ z_2$. Z własności produktu istnieje jednoznaczna strzałka $\langle z_1,z_2\rangle\colon Z\to A\times B$ z $\pi_1\circ\langle z_1,z_2\rangle = z_1$, $\pi_2\circ\langle z_1,z_2\rangle = z_2$. Widzimy, że
\[
f\circ \pi_1 \circ \langle z_1,z_2\rangle = f\circ z_1 = g\circ z_2 = g\circ \pi_2\circ\langle z_1,z_2\rangle,
\]
więc istnieje dokładnie jedna strzałka $u\colon Z\to E$ taka, że $e\circ u = \langle z_1,z_2\rangle$. Wówczas
\[
(\pi_1\circ e)\circ u = z_1, \qquad (\pi_2\circ e)\circ u = z_2.
\]
Pozostaje jednoznaczność: jeśli $u'\colon Z\to E$ również spełnia $(\pi_1\circ e)\circ u' = z_1$ oraz $(\pi_2\circ e)\circ u' = z_2$, to z jednoznaczności strzałki do produktu $e\circ u' = \langle z_1,z_2\rangle$, a stąd, z jednoznaczności w definicji ekwalizatora, $u' = u$. \qed

\section{Pullbacki, diagramy i stożki}
\subsection*{Kategorie z produktami binarnymi i ekwalizatorami a pullbacki}

\begin{stwierdzenie}
W kategorii z produktami binarnymi i ekwalizatorami dla dowolnej pary strzałek
\[
\begin{tikzcd}
A \arrow[r,"f"] & C & B \arrow[l,"g"']
\end{tikzcd}
\]
istnieje pullback. Dokładniej, biorąc:
\begin{enumerate}
  \item[(a)] $A\times B$ wraz z $\pi_1\colon A\times B \to A$, $\pi_2\colon A\times B \to B$ jako produkt,
  \item[(b)] obiekt $E$ oraz $e\colon E \to A\times B$, gdzie $E$ jest ekwalizatorem
  \[
  \begin{tikzcd}
  E \arrow[r,"e"] & A\times B \arrow[r, shift left,"f\circ\pi_1"] \arrow[r, shift right,"g\circ\pi_2"'] & C,
  \end{tikzcd}
  \]
  \item[(c)] $p_1 \stackrel{\mathrm{df}}{=} \pi_1\circ e$, $p_2 \stackrel{\mathrm{df}}{=} \pi_2\circ e$,
\end{enumerate}
otrzymujemy, że $E$ wraz ze strzałkami $p_1,p_2$ jest pullbackiem strzałek $A \xrightarrow{f} C \xleftarrow{g} B$:
\[
\begin{tikzcd}
E \arrow[r,"p_2"] \arrow[d,"p_1"'] \arrow[dr, phantom, "\lrcorner", very near start] & B \arrow[d,"g"] \\
A \arrow[r,"f"'] & C
\end{tikzcd}
\]
\end{stwierdzenie}

\paragraph{Krok 1.} Jeśli $E$ wraz z $p_1\colon E\to A$, $p_2\colon E\to B$ jest pullbackiem strzałek $A \xrightarrow{f} C \xleftarrow{g} B$, to strzałka
\[
E \xrightarrow{\langle p_1,p_2\rangle} A\times B
\]
jest ekwalizatorem $f\circ \pi_1$ i $g\circ \pi_2$:
\[
\begin{tikzcd}
E \arrow[r,"{\langle p_1,p_2\rangle}"] & A\times B \arrow[r, shift left,"f\circ\pi_1"] \arrow[r, shift right,"g\circ\pi_2"'] & C.
\end{tikzcd}
\]

\emph{Dowód.} Z założenia diagram
\[
\begin{tikzcd}
E \arrow[r,"p_2"] \arrow[d,"p_1"'] & B \arrow[d,"g"] \\
A \arrow[r,"f"'] & C
\end{tikzcd}
\]
jest przemienny. Sprawdzimy, że
\begin{enumerate}
  \item strzałka $\langle p_1,p_2\rangle$ ekwalizuje $f\circ \pi_1$ i $g\circ \pi_2$,
  \item spełnia własność uniwersalną ekwalizatora.
\end{enumerate}

\textbf{Ad 1.} Mamy $\pi_1\circ \langle p_1,p_2\rangle = p_1$ oraz $\pi_2\circ \langle p_1,p_2\rangle = p_2$, więc
\[
(f\circ\pi_1)\circ\langle p_1,p_2\rangle = f\circ p_1 = g\circ p_2 = (g\circ\pi_2)\circ\langle p_1,p_2\rangle.
\]

\textbf{Ad 2.} Weźmy $z\colon Z\to A\times B$ taki, że
\[
(f\circ\pi_1)\circ z = (g\circ\pi_2)\circ z.
\]
Oznacza to, że $f\circ(\pi_1\circ z) = g\circ(\pi_2\circ z)$, więc z własności pullbacka istnieje dokładnie jedna $u\colon Z\to E$ z
\[
p_1\circ u = \pi_1\circ z, \qquad p_2\circ u = \pi_2\circ z.
\]
Wtedy
\[
\pi_1\circ\langle p_1,p_2\rangle\circ u = p_1\circ u = \pi_1\circ z, \qquad \pi_2\circ\langle p_1,p_2\rangle\circ u = p_2\circ u = \pi_2\circ z,
\]
a stąd z jednoznaczności pary do produktu $\langle p_1,p_2\rangle\circ u = z$.

Jednoznaczność $u$ wynika z jednoznaczności w pullbacku. \qed

\bigskip

\paragraph{Krok 2.} Pokazujemy, że jeśli kategoria ma pullbacki i obiekt końcowy, to ma binarne produkty oraz ekwalizatory.

\emph{(a) Binarne produkty.} Weźmy dwa obiekty $A,B$ oraz pullback strzałek $A\xrightarrow{!_A} 1 \xleftarrow{!_B} B$:
\[
\begin{tikzcd}
P \arrow[r,"p_2"] \arrow[d,"p_1"'] \arrow[dr, phantom, "\lrcorner", very near start] & B \arrow[d,"!_B"] \\
A \arrow[r,"!_A"'] & 1
\end{tikzcd}
\]
Wtedy $P$ wraz z $p_1, p_2$ jest produktem $A\times B$.

\emph{(b) Ekwalizatory.} Weźmy parę $f, g \colon A \to B$ i rozważmy strzałki $\langle f,g\rangle\colon A\to B\times B$ oraz $\langle \id_B, \id_B\rangle\colon B\to B\times B$. Pullback
\[
\begin{tikzcd}
E \arrow[r,"u"] \arrow[d,"e"'] \arrow[dr, phantom, "\lrcorner", very near start] & B \arrow[d,"{\langle \id_B,\id_B\rangle}"] \\
A \arrow[r,"{\langle f,g\rangle}"'] & B\times B
\end{tikzcd}
\]
daje obiekt $e\colon E\to A$ będący ekwalizatorem $f$ i $g$. \qed

\begin{wniosek}
Jeśli $\C$ jest kategorią z produktami binarnymi i ekwalizatorami, to ma pullbacki.
\end{wniosek}

\begin{twierdzenie}
Kategoria ma skończone produkty i ekwalizatory wtedy i tylko wtedy, gdy ma pullbacki i obiekt końcowy.
\end{twierdzenie}

\subsection*{Diagramy i stożki}

\begin{definicja}
Niech $\Jcat$ oraz $\C$ będą kategoriami. \emph{Diagramem typu $\Jcat$ w $\C$} nazywamy dowolny funktor
\[
D\colon \Jcat\to \C.
\]
\end{definicja}

\textbf{Notacja.} W tym przypadku $\Jcat$ nazywamy kategorią indeksów. Jej obiekty oznaczamy $i,j,t,\ldots$, a jej morfizmy $\alpha,\beta,\ldots$. Zamiast $D(i)$, $D(\alpha)$ pisać będziemy $D_i$, $D_\alpha$.

\begin{definicja}
\emph{Stożkiem nad diagramem $D\colon \Jcat\to \C$} nazywamy obiekt $C\in \C$ wraz z rodziną strzałek
\[
\{c_j\colon C\to D_j\}_{j\in \Ob\Jcat}
\]
taką, że dla każdego $\alpha\colon i\to j$ w $\Jcat$ następujący diagram jest przemienny w $\C$:
\[
\begin{tikzcd}
& C \arrow[dl,"c_i"'] \arrow[dr,"c_j"] & \\
D_i \arrow[rr,"D_\alpha"'] & & D_j
\end{tikzcd}
\]
\end{definicja}

\begin{przyklad}
Weźmy kategorię $\Jcat$ generowaną przez
\[
\begin{tikzcd}
& 2 \arrow[d,"\beta"] \\
1 \arrow[r,"\alpha"'] & 3
\end{tikzcd}
\]
Jeśli rozważymy diagram $D\colon \Jcat\to \Set$, to jest on zadany przez trzy zbiory $A,B,C$ oraz 2 przekształcenia
\[
\begin{tikzcd}
& B \arrow[d,"g"] \\
A \arrow[r,"f"'] & C
\end{tikzcd}
\qquad
\text{gdzie } D_1=A,\ D_2=B,\ D_3=C,\ D_\alpha = f,\ D_\beta = g.
\]
Stożek nad tym diagramem to $S$ wraz z $s_1, s_2, s_3$ takimi, że
\[
\begin{tikzcd}
S \arrow[r,"s_2"] \arrow[d,"s_1"'] \arrow[dr,"s_3"] & B \arrow[d,"g"] \\
A \arrow[r,"f"'] & C
\end{tikzcd}
\]
jest przemienny, co jest równoważne $s_3 = f\circ s_1 = g\circ s_2$.
\end{przyklad}

\begin{przyklad}
Niech $\Jcat$ będzie kategorią dyskretną o dwóch obiektach. Diagram $D\colon \Jcat\to \C$ to po prostu para obiektów $A,B$, gdzie $D_1=A$, $D_2=B$. Stożek nad $D$ to trójka $S, s_1, s_2$:
\[
\begin{tikzcd}
& S \arrow[dl,"s_1"'] \arrow[dr,"s_2"] & \\
A & & B
\end{tikzcd}
\]
gdzie $s_1\colon S\to A$, $s_2\colon S\to B$.
\end{przyklad}

\begin{przyklad}
Niech $\Jcat$ będzie kategorią generowaną przez $1 \xrightarrow{\alpha} 2$. Diagram $D\colon \Jcat\to \C$ jest zadany przez
\[
A \xrightarrow{f} B,
\]
gdzie $D_1=A$, $D_2=B$, $D_\alpha = f$. Stożek nad $D$ to $S$ z $s_1\colon S\to A$, $s_2\colon S\to B$ takimi, że $f\circ s_1 = s_2$.
\end{przyklad}


\chapter{Granice diagramów. Kategorie kartezjańsko domknięte}
\section{Granice diagramów i ich zachowywanie}
Przypomnijmy pojęcia diagramu i stożka wprowadzone pod koniec poprzedniego rozdziału.

\begin{definicja}
Niech $\Jcat$ i $\C$ będą kategoriami. \emph{Diagramem typu $\Jcat$ w $\C$} nazywamy funktor
\[
D \colon \Jcat \longrightarrow \C.
\]
\end{definicja}

\begin{uwaga}
Kategorię $\Jcat$ nazywamy \emph{kategorią indeksów}; jej obiekty oznaczamy zwykle przez $i, j, k, \dots$, a morfizmy przez $\alpha, \beta, \gamma, \dots$. Wartości funktora $D$ oznaczamy odpowiednio przez $D_i, D_j, \dots$ oraz $D_\alpha, D_\beta, \dots$
\end{uwaga}

\begin{definicja}
\emph{Stożkiem nad diagramem $D \colon \Jcat \to \C$} nazywamy obiekt $C \in \C$ wraz z rodziną morfizmów
\[
\bigl\{\, c_i \colon C \longrightarrow D_i \,\bigr\}_{i \in \Ob \Jcat},
\]
takich, że dla każdego $\alpha \colon i \to j$ w $\Jcat$ diagram
\[
\begin{tikzcd}[row sep=large, column sep=large]
 & C \arrow[dl, "c_i"'] \arrow[dr, "c_j"] & \\
 D_i \arrow[rr, "D_\alpha"'] & & D_j
\end{tikzcd}
\]
jest przemienny.
\end{definicja}

\begin{definicja}
Niech $(C, \{c_i\})$ oraz $(C', \{c'_i\})$ będą stożkami nad diagramem $D \colon \Jcat \to \C$. \emph{Morfizmem stożków} $(C, \{c_i\}) \to (C', \{c'_i\})$ nazywamy strzałkę
\[
\vartheta \colon C \longrightarrow C' \quad \text{w } \C,
\]
taką, że diagram
\[
\begin{tikzcd}[row sep=large, column sep=large]
C \arrow[rr, "\vartheta"] \arrow[dr, "c_i"'] & & C' \arrow[dl, "c'_i"] \\
 & D_i &
\end{tikzcd}
\]
jest przemienny dla każdego $i \in \Ob \Jcat$.
\end{definicja}

\begin{definicja}
Przez $\mathrm{Cone}(D)$ oznaczamy kategorię, w której obiektami są stożki nad diagramem $D \colon \Jcat \to \C$, a morfizmami --- morfizmy stożków.
\end{definicja}

\begin{definicja}
\emph{Granicą diagramu $D \colon \Jcat \to \C$} nazywamy obiekt końcowy w kategorii $\mathrm{Cone}(D)$ i oznaczamy go przez
\[
\Bigl( \varprojlim_{j} D_j,\ \bigl\{ p_i \colon \varprojlim_{j} D_j \to D_i \bigr\} \Bigr).
\]
\end{definicja}

\subsection*{Przykłady}

\begin{przyklad}
Jeśli $\Jcat = \emptyset$, to $\mathrm{Cone}(D) = \C$, więc
\[
\varprojlim D \;=\; \text{obiekt końcowy w } \C.
\]
\end{przyklad}

\begin{przyklad}
Jeśli $\Jcat$ jest kategorią dyskretną o dwóch obiektach $\bullet_1, \bullet_2$, a $D(\bullet_i) = D_i$, to
\[
\varprojlim D \;=\; D_1 \times D_2.
\]
\end{przyklad}

\begin{przyklad}
Jeśli $\Jcat$ jest kategorią z dwoma obiektami i dwiema równoległymi strzałkami między nimi:
\[
\Jcat \;=\; \begin{tikzcd} \bullet_1 \arrow[r, shift left, "\alpha"] \arrow[r, shift right, "\beta"'] & \bullet_2 \end{tikzcd},
\qquad
D \;=\; \begin{tikzcd} D_1 \arrow[r, shift left, "D_\alpha"] \arrow[r, shift right, "D_\beta"'] & D_2 \end{tikzcd},
\]
to
\[
\varprojlim D \;=\; \text{ekwalizator } D_\alpha,\, D_\beta.
\]
\end{przyklad}

\begin{przyklad}
Jeśli $\Jcat$ jest kategorią postaci
\[
\Jcat \;=\; \begin{tikzcd} & \bullet_3 \arrow[d, "\beta"] \\ \bullet_1 \arrow[r, "\alpha"'] & \bullet_2 \end{tikzcd},
\qquad
D \;=\; \begin{tikzcd} & D_3 \arrow[d, "D_\beta"] \\ D_1 \arrow[r, "D_\alpha"'] & D_2 \end{tikzcd},
\]
to
\[
\varprojlim D \;=\; \text{pullback } D_\alpha, D_\beta \text{ w } \C.
\]
\end{przyklad}

\begin{przyklad}[Granica nieskończonego łańcucha]
Granice istnieją także nad diagramami nieskończonymi. Niech $\Jcat$ będzie kategorią indukowaną przez poset $(\mathbb{N}, \ge)$; diagram $D\colon \Jcat\to\Set$ to łańcuch zbiorów i przekształceń
\[
\cdots \longrightarrow X_2 \xrightarrow{\;f_1\;} X_1 \xrightarrow{\;f_0\;} X_0.
\]
W $\Set$ granicą takiego diagramu (zwaną \emph{granicą odwrotną}) jest
\[
\varprojlim_n X_n \;=\; \Bigl\{ (x_0, x_1, x_2, \ldots) \in \prod_n X_n \;:\; f_n(x_{n+1}) = x_n \text{ dla każdego } n \Bigr\}
\]
wraz z rzutowaniami. Na przykład dla $X_n = A^n$ (ciągi długości $n$ nad zbiorem $A$) i przekształceń obcinających ostatni element granicą jest zbiór $A^\omega$ wszystkich nieskończonych ciągów (strumieni) nad $A$.
\end{przyklad}

\begin{definicja}
\emph{Skończoną granicą} nazywamy granicę nad diagramem $D \colon \Jcat \to \C$, w którym kategoria $\Jcat$ jest skończona.
\end{definicja}

\begin{twierdzenie}
Kategoria $\C$ ma wszystkie skończone granice wtedy i tylko wtedy, gdy ma skończone produkty i ekwalizatory (lub równoważnie: gdy ma pullbacki i obiekt końcowy).
\end{twierdzenie}

\begin{proof}[Dowód (szkic)]
Pokażemy, że z faktu istnienia skończonych produktów i ekwalizatorów wynika istnienie wszystkich skończonych granic.

Weźmy diagram $D \colon \Jcat \to \C$, gdzie $\Jcat$ jest skończoną kategorią indeksów. Podamy konstrukcję obiektu $E$ wraz ze strzałkami $e_i \colon E \to D_i$, dla którego $(E, \{e_i\})$ będzie obiektem końcowym w $\mathrm{Cone}(D)$.

\textbf{(1)} Rozważmy następujące obiekty w $\C$:
\[
\prod_{i \in \Ob \Jcat} D_i \qquad \text{oraz} \qquad \prod_{\alpha \colon i \to j \in \Mor \Jcat} D_j.
\]

\textbf{(2)} Zdefiniujmy dwie strzałki
\[
\begin{tikzcd}[column sep=large]
\displaystyle\prod_{i \in \Ob \Jcat} D_i \arrow[r, shift left, "\varphi"] \arrow[r, shift right, "\psi"'] & \displaystyle\prod_{\alpha \colon i \to j} D_j,
\end{tikzcd}
\]
zadane następująco: rzutowanie $\pi_\alpha$ z produktu po prawej stronie spełnia
\[
\pi_\alpha \circ \varphi = \pi_j, \qquad \pi_\alpha \circ \psi = D_\alpha \circ \pi_i.
\]

\textbf{(3)} Weźmy ekwalizator $E$ wraz ze strzałką $e \colon E \to \prod_{i \in \Ob \Jcat} D_i$ pary $\varphi, \psi$:
\[
\begin{tikzcd}[column sep=large]
E \arrow[r, "e"] & \displaystyle\prod_{i \in \Ob \Jcat} D_i \arrow[r, shift left, "\varphi"] \arrow[r, shift right, "\psi"'] & \displaystyle\prod_{\alpha \colon i \to j} D_j.
\end{tikzcd}
\]
Definiujemy $e_i := \pi_i \circ e$. Można pokazać, że $E$ wraz z $\{e_i\}_{i \in \Ob \Jcat}$ spełnia żądane własności.
\end{proof}

\subsection*{Zachowywanie granic}

\begin{definicja}
Mówimy, że funktor $F \colon \C \to \D$ \emph{zachowuje granice} nad diagramem typu $\Jcat$, jeśli dla każdego diagramu $D \colon \Jcat \to \C$ zachodzi
\[
F\Bigl(\varprojlim_i D_i\Bigr) \;\cong\; \varprojlim_i F(D_i)
\]
(izomorfizm kanoniczny).

Mówimy, że funktor $F \colon \C \to \D$ \emph{zachowuje wszystkie granice}, jeśli zachowuje granice nad diagramami każdego typu.
\end{definicja}

\begin{przyklad}
Niech $\C$ będzie kategorią. Rozważmy funktor
\[
\C(X, -) \colon \C \longrightarrow \Set
\]
dla ustalonego obiektu $X \in \C$:
\[
A \longmapsto \C(X, A) = \text{zbiór strzałek } X \to A \text{ w } \C,
\]
oraz dla $f \colon A \to B$:
\[
\C(X, A) \xrightarrow{\;\C(X, f)\;} \C(X, B), \qquad \C(X, f)(g \colon X \to A) \;=\; f \circ g.
\]
(Często $\C(X, f) =: f_*$, czyli $f_*(g) = f \circ g$.)

Pokażemy, że funktor $\C(X, -)$ zachowuje wszystkie skończone granice.

\textbf{(1) Obiekt końcowy.} Jeśli $\mathbf{1}$ jest obiektem końcowym w $\C$, to
\[
\C(X, \mathbf{1}) \;=\; \{!_X\} \quad \text{--- jednoelementowy zbiór,}
\]
czyli obiekt końcowy w $\Set$.

\textbf{(2) Produkty.} Z ćwiczeń (zadanie~\ref{zad:hom-produkty}) wiemy, że $\C(X, -)$ zachowuje produkty.

\textbf{(3) Ekwalizatory.} Weźmy parę $f, g \colon A \rightrightarrows B$, ich ekwalizator
\[
\begin{tikzcd}
E \arrow[r, "e"] & A \arrow[r, shift left, "f"] \arrow[r, shift right, "g"'] & B.
\end{tikzcd}
\]
Czy diagram
\[
\begin{tikzcd}
\C(X, E) \arrow[r, "e_*"] & \C(X, A) \arrow[r, shift left, "f_*"] \arrow[r, shift right, "g_*"'] & \C(X, B)
\end{tikzcd}
\]
jest diagramem ekwalizatora w $\Set$? Weźmy $x \in \C(X, A)$, czyli $x \colon X \to A$. Warunek $f_*(x) = g_*(x)$ jest równoważny $f \circ x = g \circ x$. Z własności ekwalizatora istnieje dokładnie jedna strzałka $u_x \colon X \to E$ taka, że $e \circ u_x = x$. Czyli
\[
e_* \colon \C(X, E) \longrightarrow \{\, x \in \C(X, A) \colon f_*(x) = g_*(x) \,\}
\]
jest bijekcją, a więc $\C(X, E)$ jest ekwalizatorem $f_*, g_*$.
\end{przyklad}

\begin{wniosek}
Funktor $\C(X, -)$ zachowuje wszystkie skończone granice.
\end{wniosek}

\begin{uwaga}
Można pokazać, że $\C(X, -)$ zachowuje wszystkie granice (nie tylko skończone).
\end{uwaga}

\begin{przyklad}[Funktor, który nie zachowuje granic]
Nie każdy funktor zachowuje granice. Rozważmy $F = (-) + 1 \colon \Set \to \Set$, $FX = X + \{\ast\}$. Funktor $F$ nie zachowuje obiektu końcowego: $F(\mathbf{1})$ jest zbiorem dwuelementowym, a więc nie jest obiektem końcowym w $\Set$. Nie zachowuje też produktów: zbiór $F(\mathbf{1}\times\mathbf{1})$ ma dwa elementy, podczas gdy $F(\mathbf{1})\times F(\mathbf{1})$ -- cztery.
\end{przyklad}

\section{Wykładniki i kategorie kartezjańsko domknięte}
\subsection*{Wykładniki}

\begin{przyklad}
W kategorii $\Set$ mamy: dla wybranych zbiorów $B, C$ istnieje zbiór
\[
C^B \stackrel{df}{=} \{f : B \to C\}.
\]
\end{przyklad}

\begin{uwaga}
Weźmy funkcję $f : A \times B \to C$. Wybierzmy $a \in A$ i zauważmy, że
\[
f_a \stackrel{df}{=} f(a, -) : B \to C
\]
jest elementem zbioru $C^B$. Dla funkcji $f : A \times B \to C$ zdefiniowaliśmy więc funkcję
\[
\hat{f} : A \to C^B, \qquad a \mapsto f_a.
\]
Co więcej, dla każdej funkcji $g : A \to C^B$ istnieje funkcja $f : A \times B \to C$ taka, że $\hat{f} = g$ -- mianowicie $f(x, y) = g(x)(y)$.

Podsumowując, istnieje bijekcja
\[
\Set(A \times B, C) \;\cong\; \Set(A, C^B).
\]
\end{uwaga}

\begin{uwaga}
Bijekcja daje nam przekształcenie
\[
\mathrm{eval} : C^B \times B \to C, \qquad \mathrm{eval}(g, b) = g(b).
\]
Ponadto dla dowolnego $A$ i $f : A \times B \to C$ istnieje jednoznaczne $\tilde{f} : A \to C^B$ takie, że diagram
\[
\begin{tikzcd}
C^B \times B \arrow[r, "\mathrm{eval}"] & C \\
A \times B \arrow[u, "\tilde{f} \times \id_B"] \arrow[ur, "f"'] &
\end{tikzcd}
\]
jest przemienny, tj.\ $\mathrm{eval}\circ(\tilde f\times\id_B) = f$.
\end{uwaga}

\begin{definicja}
Załóżmy, że $\C$ ma binarne produkty. Obiektem potęgowym (\emph{wykładnikiem}) obiektów $B, C \in \C$ nazywamy obiekt $C^B$ wraz ze strzałką
\[
\varepsilon : C^B \times B \to C
\]
taką, że dla każdego obiektu $Z \in \C$ i strzałki $f : Z \times B \to C$ istnieje dokładnie jedna strzałka
\[
\tilde{f} : Z \to C^B
\]
taka, że diagram
\[
\begin{tikzcd}
C^B \times B \arrow[r, "\varepsilon"] & C \\
Z \times B \arrow[u, "\tilde{f} \times \id_B"] \arrow[ur, "f"'] &
\end{tikzcd}
\]
jest przemienny.
\end{definicja}

\begin{uwaga}[Notacja]
\begin{enumerate}
\item $\varepsilon : C^B \times B \to C$ nazywamy \emph{ewaluacją}.
\item $\tilde{f} : Z \to C^B$ nazywamy \emph{transpozycją} (\emph{wykładnikiem}) strzałki $f : Z \times B \to C$.
\end{enumerate}
\end{uwaga}

\begin{uwaga}
Mając daną strzałkę $g : Z \to C^B$, definiujemy
\[
\overline{g} \stackrel{df}{=} \varepsilon \circ (g \times \id_B) : Z \times B \to C,
\]
tj.
\[
\begin{tikzcd}
Z \times B \arrow[r, "g \times \id_B"] & C^B \times B \arrow[r, "\varepsilon"] & C
\end{tikzcd}
\]
Wtedy $\overline{\tilde{f}} = f$ oraz $\widetilde{\overline{g}}=g$. Innymi słowy, jeśli kategoria $\C$ ma obiekty potęgowe, to mamy bijekcję
\[
\C(Z \times B, C) \;\cong\; \C(Z, C^B), \qquad f \mapsto \tilde{f}, \quad g \mapsto \overline{g}.
\]
\end{uwaga}

\subsection*{Kategorie kartezjańsko domknięte}

\begin{definicja}
Kategorię $\C$ nazywamy \emph{kartezjańsko domkniętą}, jeśli ma:
\begin{enumerate}
\item wszystkie skończone produkty,
\item wszystkie obiekty potęgowe (dla dowolnej pary $B, C$).
\end{enumerate}
\end{definicja}

\begin{przyklad}
$\Set$ oraz $\Pos$ są kategoriami kartezjańsko domkniętymi.
\end{przyklad}

\begin{stwierdzenie}
$\Pos$ jest kartezjańsko domknięta.
\end{stwierdzenie}

\begin{proof}[Szkic]
\textbf{(1)} $\Pos$ ma wszystkie skończone produkty: jednoelementowy poset jest terminalny w $\Pos$, a dla posetów $P = (P, \leq_P)$, $Q = (Q, \leq_Q)$ definiujemy
\[
P \times Q = (P \times Q, \leq),
\]
gdzie
\[
(p_1, q_1) \leq (p_2, q_2) \;\stackrel{df}{\iff}\; p_1 \leq_P p_2 \;\wedge\; q_1 \leq_Q q_2.
\]
Razem z rzutowaniami $\pi_1, \pi_2$ daje to produkt.

\textbf{(2)} Dla posetów $P$ i $Q$ definiujemy
\[
Q^P \stackrel{df}{=} \big(\{f : P \to Q \mid f \text{ monotoniczna}\},\ \leq\big),
\]
gdzie
\[
f_1 \leq f_2 \;\stackrel{df}{\iff}\; f_1(p) \leq_Q f_2(p) \quad \text{dla każdego } p \in P,
\]
oraz ewaluację
\[
\varepsilon : Q^P \times P \to Q, \qquad \varepsilon(f, p) \stackrel{df}{=} f(p),
\]
która jest dobrze zdefiniowaną strzałką w $\Pos$. Transpozycja $\tilde{f}\colon X\to Q^P$ jest zdefiniowana analogicznie jak w $\Set$.
\end{proof}

\section{Funktor wykładniczy. Aksjomatyzacja równościowa CCC}
\begin{definicja}
Niech $\C$ będzie kategorią oraz $B, C \in \C$ obiektami. \emph{Obiektem potęgowym} (wykładniczym) nazywamy obiekt $C^B$ wraz ze strzałką
\[
\varepsilon : C^B \times B \longrightarrow C
\]
taką, że dla każdego obiektu $Z \in \C$ i każdej strzałki $f : Z \times B \to C$ istnieje dokładnie jedna strzałka $\tilde{f} : Z \to C^B$ taka, że poniższy diagram jest przemienny:
\[
\begin{tikzcd}
C^B \times B \arrow[r, "\varepsilon"] & C \\
Z \times B \arrow[u, "\tilde{f} \times \id_B"] \arrow[ur, "f"'] &
\end{tikzcd}
\]
tzn.\ $\varepsilon \circ (\tilde{f} \times \id_B) = f$.
\end{definicja}

\begin{definicja}
Kategoria $\C$ nazywana jest \emph{kartezjańsko-domkniętą} (CCC -- \emph{cartesian closed category}), gdy:
\begin{enumerate}
\item ma wszystkie skończone produkty (w szczególności obiekt końcowy $\mathbf{1}$),
\item dla każdej pary obiektów $B, C \in \C$ istnieje obiekt potęgowy $C^B$.
\end{enumerate}
\end{definicja}

\begin{stwierdzenie}
Niech $\C$ będzie CCC. Wtedy przyporządkowanie na obiektach
\[
X \longmapsto X^A
\]
rozszerza się do funktora $(-)^A : \C \to \C$.
\end{stwierdzenie}

\begin{proof}
\textbf{(1) Sytuacja w $\Set$.} Załóżmy, że mamy funkcję $\beta : B \to C$ w $\Set$. Definiujemy
\[
\beta^A : B^A \longrightarrow C^A, \qquad \beta^A(f : A \to B) \stackrel{df}{=} \beta \circ f.
\]
Mamy więc przyporządkowanie
\[
X \longmapsto X^A, \qquad (\beta : B \to C) \longmapsto (\beta^A : B^A \to C^A).
\]
Sprawdźmy aksjomaty funktorowe. Dla $\alpha : C \to D$ i $\beta : B \to C$:
\[
(\alpha \circ \beta)^A(f) = \alpha \circ \beta \circ f = (\alpha^A \circ \beta^A)(f),
\]
czyli $(\alpha \circ \beta)^A = \alpha^A \circ \beta^A$. Ponadto $(\id_B)^A(f) = \id_B \circ f = f = \id_{B^A}(f)$, więc $(\id_B)^A = \id_{B^A}$.

\textbf{(2) Ogólnie, dla CCC $\C$.} Dla strzałki $\beta : B \to C$ definiujemy $\beta^A : B^A \to C^A$ jako
\[
\beta^A \stackrel{df}{=} \widetilde{\beta \circ \varepsilon},
\]
tj.\ transpozycję strzałki $B^A \times A \xrightarrow{\varepsilon} B \xrightarrow{\beta} C$. Innymi słowy, $\beta^A$ jest jedynym rozwiązaniem równania (przemienność diagramu):
\[
\begin{tikzcd}
C^A \times A \arrow[r, "\varepsilon"] & C \\
B^A \times A \arrow[u, "\beta^A \times \id_A"] \arrow[r, "\varepsilon"'] & B \arrow[u, "\beta"']
\end{tikzcd}
\]

$(\id_B)^A$ jest jedynym rozwiązaniem
\[
\begin{tikzcd}
B^A \times A \arrow[r, "\varepsilon"] & B \\
B^A \times A \arrow[u, "(\id_B)^A \times \id_A"] \arrow[r, "\varepsilon"'] & B \arrow[u, "\id_B"']
\end{tikzcd}
\]
Ale my znamy rozwiązanie -- jest nim $\id_{B^A}$. Z jednoznaczności $(\id_B)^A = \id_{B^A}$.

Pokażemy, że $\sigma^A \circ \beta^A = (\sigma \circ \beta)^A$ dla $B \xrightarrow{\beta} C \xrightarrow{\sigma} D$. Poniższy złożony diagram jest przemienny:
\[
\begin{tikzcd}
D^A \times A \arrow[r, "\varepsilon"] & D \\
C^A \times A \arrow[u, "\sigma^A \times \id_A"] \arrow[r, "\varepsilon"] & C \arrow[u, "\sigma"'] \\
B^A \times A \arrow[u, "\beta^A \times \id_A"] \arrow[r, "\varepsilon"'] & B \arrow[u, "\beta"']
\end{tikzcd}
\]
Widzimy, że $\sigma^A \circ \beta^A$ jest rozwiązaniem równania definiującego $(\sigma \circ \beta)^A$. Z jednoznaczności:
\[
(\sigma \circ \beta)^A = \sigma^A \circ \beta^A. \qedhere
\]
\end{proof}

\begin{uwaga}
Oprócz $\varepsilon : C^B \times B \to C$ istnieje też strzałka $\eta$:
\[
\eta : C \longrightarrow (C \times B)^B
\]
zdefiniowana jako transpozycja $\id_{C \times B} : C \times B \to C \times B$. W $\Set$:
\[
\eta(c) \stackrel{df}{=} \bigl(\, b \longmapsto (c, b) \,\bigr).
\]
Jeśli przyjmiemy $F = (-)^B$ oraz $G = (-) \times B$, to mamy strzałki
\[
\eta : C \longrightarrow F G C, \qquad \varepsilon : G F C \longrightarrow C.
\]
\end{uwaga}

\begin{uwaga}
W dowolnej kategorii CCC strzałka $\eta$ pozwala znaleźć $\tilde{f}$: dla $f : Z \times A \to B$ mamy
\[
f^A : (Z \times A)^A \longrightarrow B^A,
\]
oraz
\[
\begin{tikzcd}
(Z \times A)^A \arrow[r, "f^A"] & B^A \\
Z \arrow[u, "\eta"] \arrow[ur, "\tilde{f}"'] &
\end{tikzcd}
\]
Innymi słowy $\tilde{f} = f^A \circ \eta$.
\end{uwaga}

\subsection*{Równościowa aksjomatyzacja CCC}

\begin{stwierdzenie}
Kategoria $\C$ jest CCC wtedy i tylko wtedy, gdy ma następującą strukturę:
\begin{itemize}
\item Wybrany obiekt $\mathbf{1}$ taki, że dla każdego $C \in \C$ istnieje strzałka $!_C : C \to \mathbf{1}$ oraz dla każdej $f : C \to \mathbf{1}$ zachodzi $f = \,!_C$.

\item Dla każdej pary $A, B$ obiekt $A \times B$ wraz ze strzałkami $p_1 : A \times B \to A$, $p_2 : A \times B \to B$ takimi, że dla każdej pary $f : Z \to A$, $g : Z \to B$ istnieje strzałka $\langle f, g \rangle : Z \to A \times B$ spełniająca
\[
p_1 \circ \langle f, g \rangle = f, \qquad p_2 \circ \langle f, g \rangle = g,
\]
oraz $\langle p_1 \circ h, p_2 \circ h \rangle = h$ dla każdego $h : Z \to A \times B$.

\item Dla każdej pary $A, B$ obiekt $B^A$ wraz ze strzałką $\varepsilon : B^A \times A \to B$ taką, że dla każdej $f : Z \times A \to B$ istnieje $\tilde{f} : Z \to B^A$ z $\varepsilon \circ (\tilde{f} \times \id_A) = f$ oraz dla każdego $g : Z \to B^A$ mamy $\widetilde{\varepsilon \circ (g \times \id_A)} = g$.
\end{itemize}
\end{stwierdzenie}

\section{Funktor IO i modelowanie obliczeń z efektami}
\section*{Podsumowanie najważniejszych własności CCC}

Jeśli $\C$ jest CCC, to:

\begin{enumerate}
\item Dla każdego obiektu $A$ mamy funktor
\[
(-)\times A \colon \C \longrightarrow \C, \qquad X \longmapsto X\times A, \qquad \bigl(X \xrightarrow{f} Y\bigr) \longmapsto \bigl(X\times A \xrightarrow{f\times \id_A} Y\times A\bigr).
\]

\item Dla każdego $A$ mamy funktor
\[
(-)^{A} \colon \C \longrightarrow \C, \qquad X \longmapsto X^{A},
\]
\[
\bigl(X \xrightarrow{f} Y\bigr) \longmapsto \bigl(X^{A} \xrightarrow{f^{A}} Y^{A}\bigr), \qquad f^{A} \stackrel{df}{=} \widetilde{f\circ \varepsilon}.
\]

\item Dla każdego $A,X,Y$ mamy bijekcję
\[
\C\bigl(X\times A,\; Y\bigr) \;\cong\; \C\bigl(X,\; Y^{A}\bigr),
\]
zadaną przez transpozycję $f \mapsto \tilde{f}$.

\item Dla każdego obiektu $X,A$ mamy ewaluację
\[
\varepsilon \colon X^{A}\times A \longrightarrow X.
\]
\end{enumerate}

\begin{stwierdzenie}
Dla każdego $f\colon X\to Y$ następujący diagram jest przemienny (naturalność $\varepsilon$):
\[
\begin{tikzcd}[column sep=large, row sep=large]
X^{A}\times A \arrow[r, "\varepsilon_{X}"] \arrow[d, "f^{A}\times \id_{A}"'] & X \arrow[d, "f"] \\
Y^{A}\times A \arrow[r, "\varepsilon_{Y}"'] & Y
\end{tikzcd}
\]
\end{stwierdzenie}

\begin{proof}
Weźmy
\[
g \colon X^{A}\times A \xrightarrow{\varepsilon} X \xrightarrow{f} Y.
\]
Z definicji $f^A = \widetilde{f\circ\varepsilon}$ wynika, że $\varepsilon\circ(f^A\times\id_A) = f\circ\varepsilon$, co jest dokładnie przemiennością kwadratu powyżej.
\end{proof}

\begin{enumerate}
\setcounter{enumi}{4}
\item Mamy strzałkę $\eta\colon X \to (X\times A)^{A}$ zdefiniowaną jako
\[
\eta \stackrel{df}{=} \widetilde{\id_{X\times A}}.
\]

\item Następujący wzór jest pomocny przy składaniu strzałek wykładnikowych: dla $f\colon X\times A\to Y\times A$, $g\colon Y\times A\to Z\times A$
\[
\widetilde{g\circ f} \;=\; (\varepsilon_{Z\times A})^{A}\circ \bigl(\widetilde{g}\times\id_{A}\bigr)^{A}\circ \widetilde{f}.
\]
\end{enumerate}

\begin{uwaga}
Ćwiczenie: uzasadnić powyższą równość.
\end{uwaga}

\section*{Wykorzystanie CCC w modelowaniu obliczeń z efektami ubocznymi}

Niech w naszej kategorii istnieje specjalny obiekt $W$ odpowiadający za stan ,,świata zewnętrznego''. Funkcja, która przyjmuje argument typu $X$ i zwraca wartość typu $Y$, jednocześnie ingerując w $W$, jest typu
\[
f \colon X\times W \longrightarrow Y\times W.
\]

Z bijekcji wynika
\[
\frac{f\colon X\times W \longrightarrow Y\times W}{\widetilde{f}\colon X \longrightarrow (Y\times W)^{W}}.
\]

Jeśli zdefiniujemy
\[
\IO(Y) \;=\; (Y\times W)^{W},
\]
to nasza funkcja jest typu $\widetilde{f}\colon X \to \IO(Y)$.

Tu $\IO = (-)^W \circ ((-)\times W)$ jest funktorem na $\C$.

\subsection*{Przykłady typów funkcji w Haskellu}

\[
\texttt{putChar} \;::\; \texttt{Char} \longrightarrow \IO\,()
\]
(wyświetlenie znaku na monitorze)

\[
\texttt{getChar} \;::\; \IO\,\texttt{Char}
\]

\[
\texttt{readFile} \;::\; \texttt{FilePath} \longrightarrow \IO\,\texttt{String}
\]
\[
\texttt{writeFile} \;::\; (\texttt{FilePath},\, \texttt{String}) \longrightarrow \IO\,()
\]


\chapter{Funktor stanu i kompozycja Kleisliego}\label{ch:stan}
\subsection*{Przypomnienie własności CCC}

W CCC dla dowolnych obiektów $A, B$ istnieje obiekt wykładniczy $B^A$ oraz morfizm \emph{ewaluacji}
\[
\varepsilon_{A,B} \colon B^A \times A \to B
\]
o następujących własnościach:

\textbf{(A)} Dla każdego morfizmu $f \colon X \times A \to B$ istnieje dokładnie jeden morfizm $\widetilde{f} \colon X \to B^A$ taki, że
\[
\begin{tikzcd}[column sep=large]
X \times A \arrow[r, "f"] \arrow[d, "\widetilde{f}\times \id_A"'] & B \\
B^A \times A \arrow[ur, "\varepsilon"'] &
\end{tikzcd}
\]
komutuje.

\textbf{(B)} Dla każdego $g \colon X \to B^A$ definiujemy
\[
\overline{g} \;\stackrel{df}{=}\; \varepsilon \circ (g \times \id_A) \colon X \times A \to B.
\]
Wtedy $\widetilde{\overline{g}} = g$ oraz dla każdego $f \colon X \times A \to B$ zachodzi $\overline{\widetilde{f}} = f$.

\textbf{(C)} Dla każdego $f \colon X \to Y$:
\[
\begin{tikzcd}[column sep=large]
X^A \times A \arrow[r, "\varepsilon_X"] \arrow[d, "f^A\times \id_A"'] & X \arrow[d, "f"] \\
Y^A \times A \arrow[r, "\varepsilon_Y"'] & Y
\end{tikzcd}
\]
oraz $f^A = \widetilde{f \circ \varepsilon}$.

\subsection*{Transpozycja złożenia (zadanie~\ref{zad:transpozycja})}

\begin{stwierdzenie}\label{stw:transpozycja-zlozenia}
Dla każdych
\[
f \colon X \times A \to Y \times A, \qquad g \colon Y \times A \to Z \times A
\]
zachodzi
\[
\widetilde{g \circ f} \;=\; (\varepsilon_{Z\times A})^A \circ (\widetilde{g} \times \id_A)^A \circ \widetilde{f}. \tag{$*$}
\]
\end{stwierdzenie}

\begin{proof}
Wystarczy pokazać, że wyrażenie po prawej stronie spełnia równanie definiujące $\widetilde{g \circ f}$, tzn.\ czyni diagram
\[
\begin{tikzcd}[column sep=huge, row sep=large]
X \times A \arrow[rr, "g\circ f"] \arrow[d, "(?)\times \id_A"'] & & Z \times A \\
(Z \times A)^A \times A \arrow[urr, "\varepsilon_{Z\times A}"'] & &
\end{tikzcd}
\]
przemiennym dla $(?) = (\varepsilon_{Z\times A})^A \circ (\widetilde{g} \times \id_A)^A \circ \widetilde{f}$.

Korzystając z (A) i naturalności ewaluacji (C) otrzymujemy
\[
g \circ f \;=\; \varepsilon_{Z\times A} \circ \bigl((\varepsilon_{Z\times A})^A \circ (\widetilde{g}\times\id_A)^A \circ \widetilde{f}\bigr)\times \id_A.
\]
Z jednoznaczności w (A) wynika $(*)$.
\end{proof}

\subsection*{Funktor stanu (zadanie~\ref{zad:state})}

\begin{definicja}
W kategorii CCC, mając funktory $(-)\times A$ oraz $(-)^A$, definiujemy funktor
\[
\State_A \;=\; (-)^A \circ \bigl((-)\times A\bigr),
\]
czyli
\[
\State_A X = (X\times A)^A.
\]
Działanie na strzałkach: dla $f\colon X\to Y$
\[
\State_A(f) = (f\times \id_A)^A = \widetilde{(f\times \id_A)\circ \varepsilon_{X\times A}}.
\]
\end{definicja}

\begin{definicja}[Kompozycja Kleisliego]
Dla dwóch strzałek
\[
f \colon X \to \State_A Y, \qquad g \colon Y \to \State_A Z
\]
definiujemy
\[
g \,\square\, f \;\stackrel{df}{=}\; (\varepsilon_{Z\times A})^A \circ \State_A(g) \circ f.
\]
\end{definicja}

\begin{uwaga}
Operacja $\square$ jest ,,kleislowskim'' odpowiednikiem zwykłego składania. Mamy korespondencję
\[
\begin{tikzcd}[column sep=large]
f\colon X\times A \to Y\times A \arrow[r, leftrightarrow] & \widetilde{f}\colon X \to \State_A Y
\end{tikzcd}
\]
i analogicznie dla $g$. Wówczas
\[
\frac{f\colon X\times A\to Y\times A,\quad g\colon Y\times A \to Z\times A}{g\circ f\colon X\times A \to Z\times A}
\quad\longleftrightarrow\quad
\frac{\widetilde{f},\; \widetilde{g}}{\widetilde{g}\,\square\,\widetilde{f}\colon X\to \State_A Z.}
\]
\end{uwaga}

\begin{stwierdzenie}
Operacja $\square$ jest łączna oraz spełnia
\[
\eta_Y \,\square\, f = f, \qquad f \,\square\, \eta_X = f,
\]
gdzie $\eta_X \colon X \to \State_A X$ dane jest przez $\eta_X = \widetilde{\id_{X\times A}}$.
\end{stwierdzenie}

\begin{proof}
\textbf{Łączność.} Weźmy $f \colon X \to \State_A Y$, $g \colon Y \to \State_A Z$, $h \colon Z \to \State_A T$. Z (B) mamy $f = \widetilde{\,\overline{f}\,}$, $g = \widetilde{\,\overline{g}\,}$. Wówczas
\[
g \,\square\, f \;=\; (\varepsilon_{Z\times A})^A \circ \bigl(\widetilde{\,\overline{g}\,} \times \id_A\bigr)^A \circ \widetilde{\,\overline{f}\,}.
\]
Ze Stwierdzenia~\ref{stw:transpozycja-zlozenia} prawa strona równa się $\widetilde{\,\overline{g}\circ \overline{f}\,}$, czyli
\[
g \,\square\, f \;=\; \widetilde{\,\overline{g}\circ \overline{f}\,}. \tag{$\triangle$}
\]
Dalej
\begin{align*}
h \,\square\, (g \,\square\, f) &= h \,\square\, \widetilde{\,\overline{g}\circ \overline{f}\,} \stackrel{(\triangle)}{=} \widetilde{\,\overline{h}\circ (\overline{g}\circ \overline{f})\,} \\
 &= \widetilde{\,(\overline{h}\circ \overline{g})\circ \overline{f}\,} \stackrel{(\triangle)}{=} (h \,\square\, g) \,\square\, f.
\end{align*}

\textbf{Elementy neutralne.} Mamy $\overline{\eta_Y} = \id_{Y\times A}$ (z (B)). Z $(\triangle)$:
\[
\eta_Y \,\square\, f \;=\; \widetilde{\,\overline{\eta_Y}\circ \overline{f}\,} \;=\; \widetilde{\,\overline{f}\,} \;=\; f.
\]
Analogicznie $f \,\square\, \eta_X = f$. \qedhere
\end{proof}


\chapter{Monady}\label{ch:monady}
Niech $\C$ będzie kategorią kartezjańsko zamkniętą (CCC). Przypomnijmy zwykłe składanie morfizmów:
\[
\frac{X \xrightarrow{f} Y \qquad Y \xrightarrow{g} Z}{X \xrightarrow{g \circ f} Z}
\]

\begin{uwaga}[Motywacja --- ,,obliczenia z efektem'']
Ustalmy obiekt $A \in \Ob(\C)$ i oznaczmy
\[
\IO X \;\stackrel{df}{=}\; (X \times A)^A.
\]
Chcielibyśmy mieć regułę składania morfizmów ,,z efektem'' postaci
\[
\frac{X \xrightarrow{f} \IO Y \qquad Y \xrightarrow{g} \IO Z}{X \xrightarrow{g \,\square\, f} \IO Z}
\]
oraz wyróżnioną rodzinę morfizmów
\[
\eta = \{\eta_X : X \to \IO X\}_{X \in \C}.
\]
\end{uwaga}

\begin{przyklad}[Motywacja --- Haskell]
Napisać program, który:
\begin{enumerate}
\item wczytuje z klawiatury nazwę pliku,
\item usuwa z zawartości wielkie litery,
\item wyświetla zawartość wypisaną od tyłu.
\end{enumerate}

Dane są funkcje
\begin{align*}
\mathrm{readString} &: () \to \IO\,\mathrm{String}, \\
\mathrm{putString} &: \mathrm{String} \to \IO\,(), \\
\mathrm{readFile} &: \mathrm{String} \to \IO\,\mathrm{String}, \\
\mathrm{reverse} &: \mathrm{String} \to \mathrm{String}, \\
\mathrm{removeCaps} &: \mathrm{String} \to \mathrm{String},
\end{align*}
oraz operacje $\eta$ (czyli \texttt{pure} $:\ a \to \IO a$) i $\square$. Wówczas
\[
\mathrm{main} \;=\; \mathrm{putString} \;\square\; (\mathrm{pure} \circ \mathrm{reverse}) \;\square\; (\mathrm{pure} \circ \mathrm{removeCaps}) \;\square\; \mathrm{readFile} \;\square\; \mathrm{readString}.
\]
\end{przyklad}

\begin{uwaga}[Pytanie]
Jakie własności abstrakcyjne pozwalają na zdefiniowanie ,,składania''
\[
\frac{X \xrightarrow{f} \IO Y \qquad Y \xrightarrow{g} \IO Z}{X \xrightarrow{g \,\square\, f} \IO Z} \;?
\]

\textbf{Odpowiedź} (por.\ zad.\ 4.7): dane są
\[
\eta = \{\eta_X : X \to \IO X\}_{X},
\qquad
\mu = \bigl\{(\varepsilon_{X \times A})^A : \IO^2 X \to \IO X\bigr\}_{X},
\]
gdzie $\varepsilon$ jest kojednością sprzężenia $(-)\times A \dashv (-)^A$.
\end{uwaga}

\begin{uwaga}[Własności rodzin $\eta$ i $\mu$]
\begin{enumerate}
\item $\eta$ jest naturalną transformacją $\Id \Rightarrow \IO$ -- dla każdego $f : X \to Y$:
\[
\begin{tikzcd}
X \arrow[r, "\eta_X"] \arrow[d, "f"'] & \IO X \arrow[d, "\IO f"] \\
Y \arrow[r, "\eta_Y"'] & \IO Y
\end{tikzcd}
\]
komutuje.
\item Analogicznie $\mu$ jest naturalną transformacją $\IO^2 \Rightarrow \IO$.
\item Mamy
\[
\eta_X \;=\; \widetilde{\id_{X \times A}} \;:\; X \to (X \times A)^A \;=\; \IO X.
\]
\end{enumerate}
\end{uwaga}

\begin{uwaga}[Aksjomaty jedności i łączności]
Rodziny $\eta$ i $\mu$ spełniają dodatkowo:
\[
\begin{tikzcd}[column sep=large]
\IO X \arrow[r, "\eta_{\IO X}"] \arrow[dr, "\id_{\IO X}"'] & \IO^2 X \arrow[d, "\mu_X"] & \IO X \arrow[l, "\IO(\eta_X)"'] \arrow[dl, "\id_{\IO X}"] \\
& \IO X &
\end{tikzcd}
\qquad
\begin{tikzcd}
\IO^3 X \arrow[r, "\mu_{\IO X}"] \arrow[d, "\IO(\mu_X)"'] & \IO^2 X \arrow[d, "\mu_X"] \\
\IO^2 X \arrow[r, "\mu_X"'] & \IO X
\end{tikzcd}
\]
gdzie operacja $\square$ jest zdefiniowana wzorem
\[
g \,\square\, f \;=\; \mu_Z \circ \IO(g) \circ f, \qquad f : X \to \IO Y,\ g : Y \to \IO Z.
\]
\end{uwaga}

\section{Definicja monady}

\begin{definicja}
\textbf{Monadą} w kategorii $\C$ nazywamy trójkę $(T, \mu, \eta)$, gdzie
\begin{enumerate}
\item $T : \C \to \C$ jest funktorem,
\item $\mu = \{\mu_X : T^2 X \to T X\}_X$ jest naturalną transformacją $\mu : T^2 \Rightarrow T$,
\item $\eta = \{\eta_X : X \to T X\}_X$ jest naturalną transformacją $\eta : \Id_{\C} \Rightarrow T$,
\end{enumerate}
oraz następujące diagramy komutują:
\[
\begin{tikzcd}[column sep=large]
T X \arrow[r, "\eta_{T X}"] \arrow[dr, "\id_{T X}"'] & T^2 X \arrow[d, "\mu_X"] & T X \arrow[l, "T(\eta_X)"'] \arrow[dl, "\id_{T X}"] \\
& T X &
\end{tikzcd}
\qquad
\begin{tikzcd}
T^3 X \arrow[r, "\mu_{T X}"] \arrow[d, "T(\mu_X)"'] & T^2 X \arrow[d, "\mu_X"] \\
T^2 X \arrow[r, "\mu_X"'] & T X
\end{tikzcd}
\]
(odpowiednio: \emph{aksjomat jedności} oraz \emph{aksjomat łączności}).
\end{definicja}

\begin{uwaga}
Nieformalnie: jeśli wyobrazimy sobie $A$ jako punkt, $TA$ jako ,,$A$ owinięte pojedynczą warstwą'', a $T^2 A$ jako ,,$A$ owinięte podwójną warstwą'', to $\eta_A : A \to TA$ pakuje punkt w warstwę, a $\mu_A : T^2 A \to TA$ scala dwie warstwy w jedną. Aksjomat jedności mówi, że doklejenie warstwy a następnie scalanie nic nie zmienia; aksjomat łączności mówi, że scalanie trzech warstw nie zależy od kolejności.
\end{uwaga}

\begin{stwierdzenie}
Niech $(T, \mu, \eta)$ będzie monadą w $\C$. Wówczas operacja
\[
g \,\square\, f \;=\; \mu_Z \circ T(g) \circ f \qquad \text{dla } f : X \to T Y,\ g : Y \to T Z
\]
jest
\begin{enumerate}
\item łączna: $h \,\square\, (g \,\square\, f) = (h \,\square\, g) \,\square\, f$,
\item ze stronami neutralnymi: $\eta_Y \,\square\, f = f = f \,\square\, \eta_X$.
\end{enumerate}
\end{stwierdzenie}

\begin{uwaga}
Dowód na ćwiczeniach.
\end{uwaga}


\chapter{Kategorie Kleisliego}\label{ch:kleisli}
W poprzednim rozdziale poznaliśmy pojęcie monady $(T,\mu,\eta)$ oraz pomocniczą operację składania
\[
g\,\square\,f \;=\; \mu_Z\circ T(g)\circ f, \qquad f\colon X\to TY,\ g\colon Y\to TZ.
\]
Tę operację można teraz wprost przekuć na strukturę nowej kategorii.

\section{Definicja kategorii Kleisliego}

\begin{definicja}
Niech $(T,\mu,\eta)$ będzie monadą w kategorii $\C$. \emph{Kategorią Kleisliego} $\Kl(T)$ nazywamy kategorię, w której:
\begin{enumerate}
\item obiektami $\Kl(T)$ są obiekty $\C$;
\item dla obiektów $X,Y\in\C$ zbiór strzałek $X\to Y$ w $\Kl(T)$ definiujemy jako
\[
\Kl(T)(X,Y) \;\stackrel{df}{=}\; \C(X,\,T Y);
\]
\item identycznością na obiekcie $X$ jest jedność monady $\eta_X\colon X\to TX$;
\item dla strzałek $f\colon X\to Y$ i $g\colon Y\to Z$ w $\Kl(T)$ (tzn.\ $f\colon X\to TY$, $g\colon Y\to TZ$ w $\C$) złożenie $g\circ_{\Kl} f\colon X\to Z$ w $\Kl(T)$ dane jest wzorem
\[
g\circ_{\Kl} f \;\stackrel{df}{=}\; \mu_Z \circ T(g) \circ f \;=\; g\,\square\,f.
\]
\end{enumerate}
\end{definicja}

Schemat składania:
\[
\begin{tikzcd}[column sep=large]
X \arrow[r,"f"] & TY \arrow[r,"T(g)"] & T^2 Z \arrow[r,"\mu_Z"] & TZ.
\end{tikzcd}
\]

\begin{stwierdzenie}
$\Kl(T)$ jest kategorią.
\end{stwierdzenie}

\begin{proof}[Szkic]
Sprawdzamy aksjomaty kategorii. Niech $f\colon X\to TY$, $g\colon Y\to TZ$, $h\colon Z\to TW$ w $\C$.

\emph{Łączność.} Korzystając z naturalności $\mu$ oraz z aksjomatu łączności $\mu_W\circ\mu_{TW} = \mu_W\circ T(\mu_W)$ liczymy
\begin{align*}
h\circ_{\Kl}(g\circ_{\Kl} f) &= \mu_W \circ T(h) \circ \bigl(\mu_Z\circ T(g)\circ f\bigr) \\
 &= \mu_W \circ T(h)\circ \mu_Z \circ T(g)\circ f \\
 &\stackrel{\mathrm{nat.}\,\mu}{=} \mu_W \circ \mu_{TW} \circ T^2(h) \circ T(g)\circ f \\
 &\stackrel{\mathrm{ass.}\,\mu}{=} \mu_W \circ T(\mu_W) \circ T^2(h) \circ T(g)\circ f \\
 &= \mu_W \circ T\bigl(\mu_W \circ T(h)\circ g\bigr) \circ f \\
 &= (h\circ_{\Kl} g)\circ_{\Kl} f.
\end{align*}

\emph{Prawa jedynka.} Niech $f\colon X\to Y$ w $\Kl(T)$, tzn.\ $f\colon X\to TY$ w $\C$. Z aksjomatu jedności monady $\mu_Y \circ T(\eta_Y) = \id_{TY}$ otrzymujemy
\[
\eta_Y \circ_{\Kl} f \;=\; \mu_Y \circ T(\eta_Y) \circ f \;=\; \id_{TY}\circ f \;=\; f.
\]

\emph{Lewa jedynka.} Z naturalności $\eta\colon \Id\Rightarrow T$ mamy $T(f)\circ \eta_X = \eta_{TY}\circ f$. Stąd
\[
f \circ_{\Kl} \eta_X \;=\; \mu_Y \circ T(f) \circ \eta_X \;=\; \mu_Y \circ \eta_{TY}\circ f \;\stackrel{\eta\text{-lewa}}{=}\; \id_{TY}\circ f \;=\; f. \qedhere
\]
\end{proof}

\section{Przykłady}

\begin{przyklad}[Funkcje częściowe]
Niech $T = \mathrm{Maybe}\colon \Set\to\Set$, $TX = X+\{\bot\}$, $\eta_X(x)=x$, $\mu_X\colon (X+\{\bot\})+\{\bot\}\to X+\{\bot\}$ zadane przez sklejenie obu kopii $\bot$ w jeden. Wówczas
\[
\Kl(\mathrm{Maybe})(X,Y) = \Set(X, Y+\{\bot\}),
\]
czyli strzałki $X\to Y$ w $\Kl(\mathrm{Maybe})$ to dokładnie funkcje częściowe z $X$ do $Y$ (z $\bot$ jako wartością nieokreśloną). Składanie wyznaczone przez Kleisliego pokrywa się ze składaniem funkcji częściowych z rozdziału~1. Otrzymujemy
\[
\Kl(\mathrm{Maybe}) \;\cong\; \Par.
\]
\end{przyklad}

\begin{przyklad}[Trace nad monoidem]
Niech $(M,\cdot,1)$ będzie monoidem i niech $T = M\times(-)\colon \Set\to\Set$. Definiujemy
\[
\eta_X\colon X\to M\times X, \qquad \eta_X(x) = (1,x),
\]
\[
\mu_X\colon M\times(M\times X)\to M\times X, \qquad \mu_X(m_1,(m_2,x)) = (m_1\cdot m_2,\, x).
\]
\begin{uwaga}
Aksjomaty jedności i łączności wynikają wprost z aksjomatów monoidu (jedność jedyny po obu stronach, łączność mnożenia).
\end{uwaga}
Wówczas
\[
\Kl(M\times{-})(X,Y) \;=\; \Set\bigl(X,\, M\times Y\bigr),
\]
a złożenie strzałek $f\colon X\to M\times Y$, $g\colon Y\to M\times Z$ dane jest przez
\[
(g\,\square\, f)(x) = (m_f\cdot m_g,\, z), \qquad \text{gdzie } f(x)=(m_f,y),\ g(y)=(m_g,z).
\]
Jest to dokładnie kategoria $\Trace$ wprowadzona w rozdziale~1.
\end{przyklad}

\begin{przyklad}[Niedeterminizm i relacje]
Niech $T = \mathcal{P}\colon \Set\to\Set$ -- funktor potęgowy. Z jednościami $\eta_X(x)=\{x\}$ oraz mnożeniem $\mu_X(\mathcal{A}) = \bigcup\mathcal{A}$ funktor $\mathcal{P}$ tworzy monadę. Strzałka $X\to Y$ w $\Kl(\mathcal{P})$ to funkcja $X\to \mathcal{P}Y$, czyli równoważnie relacja $X\to Y$. Składanie Kleisliego pokrywa się ze składaniem relacji:
\[
(g\,\square\, f)(x) \;=\; \bigcup_{y\in f(x)} g(y).
\]
Otrzymujemy więc izomorfizm kategorii
\[
\Kl(\mathcal{P}) \;\cong\; \Rel.
\]
\end{przyklad}

\begin{przyklad}[Monada stanu]
Niech $A\in\Ob(\C)$ i niech $\C$ będzie kategorią kartezjańsko domkniętą. Z rozdziału~12 wiemy, że
\[
\State_A \;=\; (-\times A)^A\colon \C\to\C
\]
wraz z $\eta_X = \widetilde{\id_{X\times A}}$ oraz $\mu_X = (\varepsilon_{X\times A})^A$ jest monadą. Wówczas
\[
\Kl(\State_A)(X,Y) \;=\; \C\bigl(X,\, (Y\times A)^A\bigr) \;\cong\; \C\bigl(X\times A,\, Y\times A\bigr).
\]
Stąd intuicja: strzałka w $\Kl(\State_A)$ to ,,obliczenie z efektem ubocznym na stanie typu $A$'', tj.\ procedura, która biorąc wejście typu $X$ i obecny stan, produkuje wynik typu $Y$ i nowy stan.

Operacja $\square$ z rozdziału~11 jest dokładnie składaniem w $\Kl(\State_A)$.
\end{przyklad}

\section{Sprzężenie $F\dashv G$}

Kategoria Kleisliego jest powiązana z $\C$ przez parę naturalnych funktorów.

\begin{definicja}
Definiujemy funktor
\[
F\colon \C\longrightarrow \Kl(T)
\]
wzorami
\[
F(X) = X, \qquad F(f\colon X\to Y) = \eta_Y\circ f\colon X\to TY,
\]
oraz funktor
\[
G\colon \Kl(T)\longrightarrow \C
\]
wzorami
\[
G(X) = TX, \qquad G(f\colon X\to Y \in \Kl(T)) = \mu_Y\circ T(f)\colon TX\to TY,
\]
gdzie $f\colon X\to TY$ traktujemy jako strzałkę w $\C$.
\end{definicja}

\begin{uwaga}
Można sprawdzić bezpośrednio, że $F$ i $G$ są funktorami oraz że
\[
G\circ F \;=\; T.
\]
Co więcej, dla dowolnych $X\in\C$ oraz $Y\in\Kl(T)$ zachodzi bijekcja
\[
\Kl(T)\bigl(F(X),\, Y\bigr) \;=\; \Kl(T)(X,Y) \;=\; \C(X,TY) \;=\; \C\bigl(X,\, G(Y)\bigr),
\]
naturalna względem $X$ i $Y$. Oznacza to, że $F$ jest \emph{lewym sprzężonym} do $G$, w notacji $F\dashv G$, a wyjściowa monada $(T,\mu,\eta)$ jest dokładnie monadą indukowaną przez to sprzężenie ($T = G\circ F$, $\eta$ -- jedność, $\mu = G\varepsilon F$, gdzie $\varepsilon$ jest kojednością).
\end{uwaga}

\begin{uwaga}
Konstrukcja Kleisliego ma własność uniwersalną: spośród wszystkich par funktorów $\C\xrightarrow{F'} \Dcat \xrightarrow{G'} \C$ generujących daną monadę $T$, $\Kl(T)$ jest \emph{początkowa} (równoważnie: każdą taką parę $(F',G')$ da się jednoznacznie rozłożyć przez $\Kl(T)$). Rolą dualną pełni kategoria algebr Eilenberga--Moore'a $T\text{-}\catname{Alg}$, terminalna wśród rozkładów monady.
\end{uwaga}

\section{Podsumowanie}

Powtórzmy najważniejsze fakty:

\begin{itemize}[leftmargin=*]
\item Każda monada $(T,\mu,\eta)$ w kategorii $\C$ wyznacza kategorię $\Kl(T)$, której strzałkami są \emph{obliczenia z efektem opisywanym przez $T$}.
\item Identyczność w $\Kl(T)$ to jedność monady $\eta$; składanie -- operacja $\square$ wprowadzona w rozdziale~11.
\item Konkretne przykłady ilustrują, że pojęcie monady jednolicie obejmuje funkcje częściowe ($\mathrm{Maybe}$), efekty ,,trace'' ($M\times{-}$), niedeterminizm ($\mathcal{P}$) oraz stan ($\State_A$).
\item Konstrukcja $\Kl(T)$ wpisuje monadę $T$ w naturalne sprzężenie $F\dashv G$ między $\C$ a $\Kl(T)$.
\end{itemize}


\renewcommand{\thechapter}{Z}
\chapter{Zadania}
Poniższe zadania pochodzą z pięciu serii ćwiczeń towarzyszących wykładom z teorii kategorii. Numeracja sekcji odpowiada oryginalnej numeracji serii.

\section{Kategorie i funktory}

\begin{zadanie}
Niech $A$ będzie zbiorem. Pokazać, że przyporządkowania $\Set\to\Set$ zdefiniowane na obiektach i morfizmach jak poniżej są funktorami:
\begin{itemize}
\item $X\mapsto A\times X$ oraz $(f\colon X\to Y)\mapsto \bigl((\id\times f)\colon A\times X\to A\times Y;\, (a,x)\mapsto (a,f(x))\bigr)$,
\item $X\mapsto A+X$ oraz
\[
(f\colon X\to Y)\mapsto (\id+f)\colon A+X\to A+Y;\; x\mapsto
\begin{cases} f(x) & \text{jeśli } x\in X, \\ x & \text{jeśli } x\in A. \end{cases}
\]
\item $X\mapsto X^A$ oraz
\[
(f\colon X\to Y)\mapsto f^A\colon X^A\to Y^A;\; \varphi\mapsto f\circ\varphi.
\]
\item $X\mapsto \mathcal{P}X \stackrel{df}{=}\{A\subseteq X\}$ oraz
\[
(f\colon X\to Y)\mapsto \mathcal{P}(f)\colon \mathcal{P}(X)\to\mathcal{P}(Y);\; A\mapsto f(A).
\]
\end{itemize}
\end{zadanie}

\begin{zadanie}
Niech $\mathbb{P}$ i $\mathbb{Q}$ będą kategoriami zadanymi przez posety $(P,\leq)$ i $(Q,\leq)$. Podać charakteryzację funktorów $\mathbb{P}\to\mathbb{Q}$ wyrażoną w języku przekształceń $P\to Q$ między tymi posetami.
\end{zadanie}

\begin{zadanie}
Podać przykłady funktorów $\Set\to\Par$ i $\Par\to\Set$.
\end{zadanie}

\begin{zadanie}
Pokazać, że funktor $(-)^*\colon \Set\to\Set$ spełnia
\[
X^* \;\cong\; \{\varepsilon\} + X\times X^*,
\]
gdzie $\cong$ oznacza bijekcję między zbiorami. Dodatkowo pokazać, że dla $f\colon X\to Y\in\Set$ przekształcenie $f^*\colon X^*\to Y^*$ spełnia
\[
f^*(\varepsilon)=\varepsilon, \qquad f^*(x::xs) = f(x)::f^*(xs),
\]
gdzie operacja $({::})\colon X\times X^*\to X^*$ jest dopisaniem elementu na początku ciągu.
\end{zadanie}

\begin{zadanie}
Znaleźć przykład funktora $F\colon\Set\to\Set$ spełniającego
\[
F X \;\cong\; \{\mathrm{Nil}\} + F X \times X \times F X.
\]
\end{zadanie}

\begin{zadanie}
Pokazać, że $\Rel$ jest izomorficzna z $\Rel^{\op}$.
\end{zadanie}

\begin{zadanie}
Pokazać, że kategoria $\Set$ nie jest izomorficzna z kategorią $\Set^{\op}$.
\end{zadanie}

\begin{zadanie}
Pokazać, że w kategorii $\Set$ strzałka $f\colon A\to B$ jest izomorfizmem wtedy i tylko wtedy, gdy jest bijekcją.
\end{zadanie}

\begin{zadanie}
Pokazać, że każdy izomorfizm jest mono i epi.
\end{zadanie}

\begin{zadanie}
Pokazać, że w $\Set$ przekształcenie $f\colon X\to Y$ jest ,,na'' wtedy i tylko wtedy, gdy jest epimorfizmem.
\end{zadanie}

\begin{zadanie}
Niech $f\colon X\to Y$, $g\colon Y\to Z$ oraz $h\colon X\to Z$ spełniają $h=g\circ f$. Pokazać, że:
\begin{itemize}
\item jeśli $f,g$ są izo, to $h$ też jest izo;
\item jeśli $h$ jest mono, to $f$ jest mono;
\item jeśli $h$ jest mono, to $g$ \emph{nie} musi być mono.
\end{itemize}
\end{zadanie}

\section{Podstawowe konstrukcje kategoryjne}

\begin{zadanie}
Znaleźć obiekt początkowy i końcowy (jeśli istnieją) w kategorii $\Grp$, $\Pos$, $\Par$ oraz w kategorii wszystkich algebr typu $(1,0)$.
\end{zadanie}

\begin{zadanie}
Niech $\C$ będzie kategorią oraz niech $A\in\C$ będzie obiektem. Definiujemy przyporządkowanie $\C(A,-)\colon\C\to\Set$ wzorami
\[
X\mapsto \C(A,X), \qquad (f\colon X\to Y)\mapsto \bigl(\C(A,f)\colon\C(A,X)\to\C(A,Y);\, g\mapsto f\circ g\bigr).
\]
Pokazać, że $\C(A,-)$ jest funktorem. Udowodnić, że zachowuje binarne produkty.\footnote{Mówimy, że funktor $F\colon\C\to\D$ \emph{zachowuje produkty}, jeśli $F(A\times B)\cong F(A)\times F(B)$ dla dowolnych $A,B\in\C$, dla których produkt $A\times B$ istnieje.}
\end{zadanie}

\begin{zadanie}
Pokazać, że w dowolnej kategorii $\C$ z (binarnymi) produktami dla dowolnych trzech obiektów $A,B,C$ zachodzi
\[
(A\times B)\times C \;\cong\; A\times(B\times C).
\]
\end{zadanie}

\begin{zadanie}
Dla dowolnej rodziny $(X_i)_{i\in I}$ obiektów w kategorii $\C$ napisać definicję produktu $\prod_{i\in I} X_i$ przez własność uniwersalności, uogólniając przypadek $|I|=2$.
\end{zadanie}

\begin{zadanie}
Pokazać, że binarne koprodukty są wyznaczone jednoznacznie z dokładnością do izomorfizmu.
\end{zadanie}

\begin{zadanie}
Udowodnić, że w kategorii grup abelowych $\Ab$ koprodukt $A+B$ dwóch grup spełnia $A+B\cong A\times B$.
\end{zadanie}

\begin{zadanie}
Pokazać, że dla dwóch zbiorów $A$, $B$ koprodukt monoidów słów $A^*$ oraz $B^*$ w kategorii $\Mon$ istnieje i spełnia
\[
A^* + B^* \;\cong\; (A+B)^*.
\]
\end{zadanie}

\begin{zadanie}
Pokazać, że jeśli $e\colon E\to A$ jest ekwalizatorem pewnej pary strzałek, to $e$ jest mono. Sformułować i udowodnić (bez używania zasady dualności) twierdzenie dualne dla koekwalizatorów.
\end{zadanie}

\begin{zadanie}
Pokazać, że kategoria grup abelowych $\Ab$ ma wszystkie ekwalizatory. Opisać je.
\end{zadanie}

\begin{zadanie}
Pokazać, że $\Set$ ma wszystkie koekwalizatory. Podać ich konstrukcję.
\end{zadanie}

\begin{zadanie}
Niech $f\colon A\to B$ będzie pewnym przekształceniem w $\Set$. Opisać ekwalizator przekształceń $f\circ\pi_1,\, f\circ\pi_2\colon A\times A\to B$ (taki typ ekwalizatora nazywać będziemy \emph{jądrem} $f$). Udowodnić, że dla dowolnej relacji równoważności $R$ na $A$ jądrem przekształcenia ilorazowego $A\to A/R;\, a\mapsto[a]_R$ jest $R$.
\end{zadanie}

\begin{zadanie}
Niech $R\subseteq A\times A$ będzie dowolną relacją binarną oraz niech $\langle R\rangle$ będzie najmniejszą relacją równoważności na $A$ zawierającą $R$. Pokazać, że przekształcenie ilorazowe $A\to A/\langle R\rangle$ jest koekwalizatorem dwóch rzutowań $R\rightrightarrows A$.
\end{zadanie}

\section{(Ko)granice}

\begin{zadanie}
W kategorii $\Set$ rozważmy dowolne przekształcenie $f\colon A\to B$ oraz zanurzenie $i_U\colon U\hookrightarrow B$. Opisać pullback przekształceń $f$ i $i_U$.
\end{zadanie}

\begin{zadanie}
Pokazać, że strzałka $m\colon X\to Y$ jest mono wtedy i tylko wtedy, gdy $X$ wraz z parą $\id\colon X\to X$, $\id\colon X\to X$ jest pullbackiem pary $m\colon X\to Y$, $m\colon X\to Y$. Wywnioskować, że funktor $\C(A,-)$ zachowuje monomorfizmy.
\end{zadanie}

\begin{zadanie}
Pokazać, że w dowolnej kategorii dla pullbacku $A'\xleftarrow{m'} M' \to M$ morfizmów $A'\xrightarrow{f} A\xleftarrow{m} M$, jeśli $m$ jest mono, to $m'$ też jest mono.
\end{zadanie}

\begin{zadanie}
Pokazać, że jeśli $\C$ jest kategorią z wszystkimi skończonymi granicami oraz $F\colon\C\to\C$ jest funktorem, to kategoria $\mathrm{Alg}(F)$ ma wszystkie skończone granice.
\end{zadanie}

\begin{zadanie}
Pokazać, że jeśli kategoria $\C$ ma skończone produkty i pullbacki, to ma ekwalizatory.
\end{zadanie}

\begin{zadanie}
Zdualizować definicję pullbacku (nazywając go \emph{pushoutem}). Opisać pushouty w $\Set$. Pokazać konstrukcję pushoutów za pomocą koproduktów i koekwalizatorów.
\end{zadanie}

\begin{zadanie}
Podać definicję kostożka i kogranicy nad diagramem $D\colon\Jcat\to\mathcal{K}$. Pokazać, że kategoria $\mathcal{K}$ ma wszystkie skończone kogranice wtedy i tylko wtedy, gdy ma wszystkie skończone koprodukty i koekwalizatory (bez używania zasady dualności).
\end{zadanie}

\begin{zadanie}
Pokazać, że funktor $\C(-,A)\colon\C^{\op}\to\Set$ przekształca koekwalizatory w $\C$ na ekwalizatory, a obiekt początkowy w $\C$ na obiekt końcowy w $\Set$.
\end{zadanie}

\section{Kategorie kartezjańsko domknięte}

\begin{zadanie}
Niech $\C$ będzie CCC. Pokazać, że $\tilde{f} = f^A\circ\eta$, gdzie dla $f\colon Z\times A\to B$ strzałka $\tilde{f}\colon Z\to B^A$ oznacza transpozycję, $f^A\colon (Z\times A)^A\to B^A$ oraz $\eta\colon Z\to (Z\times A)^A$ są zdefiniowane jak na wykładzie.
\end{zadanie}

\begin{zadanie}
Pokazać, że w dowolnej kategorii kartezjańsko domkniętej zachodzi:
\begin{itemize}
\item $(A\times B)^C \cong A^C\times B^C$,
\item $(A^B)^C \cong A^{B\times C}$.
\end{itemize}
\end{zadanie}

\begin{zadanie}
Czy kategoria $\Mon$ jest CCC?
\end{zadanie}

\begin{zadanie}
Pokazać, że kategoria $\omega\mathrm{CPO}$ jest CCC, natomiast kategoria $\omega\mathrm{CPO}_\bot$ nie jest CCC.\footnote{Poset $(P,\leq)$ nazywamy $\omega\mathrm{CPO}$, jeśli każdy przeliczalny łańcuch $x_1\leq x_2\leq\dots$ ma supremum. Przekształcenie $f\colon P\to Q$ zachowujące porządek między $\omega\mathrm{CPO}$ nazywamy \emph{ciągłym}, jeśli zachowuje suprema przeliczalnych łańcuchów: $f(\bigvee_i x_i) = \bigvee_i f(x_i)$. $\omega\mathrm{CPO}$ wraz z ciągłymi przekształceniami tworzą kategorię $\omega\mathrm{CPO}$. Punktowe $\omega\mathrm{CPO}$ (z elementem najmniejszym $\bot$) wraz z ciągłymi przekształceniami zachowującymi $\bot$ tworzą kategorię $\omega\mathrm{CPO}_\bot$.}
\end{zadanie}

\begin{zadanie}
Pokazać, że kategoria $\mathbf{Cat}$ wszystkich małych kategorii i funktorów jest CCC, gdzie $\C^\D = \mathrm{Fun}(\D,\C)$.
\end{zadanie}

\begin{zadanie}
Udowodnić (dokończając dowód z wykładu), że dla każdego $f\colon X\times A\to Y\times A$ i $g\colon Y\times A\to Z\times A$ zachodzi
\[
\widetilde{g\circ f} \;=\; (\varepsilon_{Z\times A})^A \circ (\widetilde{g}\times\id_A)^A \circ \widetilde{f}.
\]
\end{zadanie}

\begin{zadanie}
Niech $\State_A$ będzie złożeniem funktorów $(-)\times A$ i $(-)^A$, tj.
\[
\State_A(X) = (X\times A)^A, \qquad \State_A(X\xrightarrow{f} Y) = \widetilde{(f\times\id_A)\circ\varepsilon_{X\times A}}.
\]
Dla strzałek $f\colon X\to\State_A Y$, $g\colon Y\to\State_A Z$ definiujemy
\[
g\cdot f \;=\; (\varepsilon_{Z\times A})^A \circ \State_A(g) \circ f.
\]
Pokazać, że tak zdefiniowane działanie jest łączne oraz $\eta_Y\cdot f = f\cdot \eta_X = f$ dla $\eta_X = \widetilde{\id_{X\times A}}$.

\emph{Wskazówka:} skorzystać z poprzedniego zadania.
\end{zadanie}

\begin{zadanie}
Niech $W$ będzie wybranym obiektem kategorii $\C$ i przyjmijmy $\IO = \State_W$. Używając operacji składania $\circ$ w $\C$, operacji $\cdot$ zdefiniowanej powyżej oraz morfizmów
\begin{align*}
& \mathbf{show}\colon\mathrm{Bool}\to\mathrm{String},\quad \mathbf{length}\colon\mathrm{String}\to\mathrm{Integer},\\
& \mathbf{removeCapitals}\colon\mathrm{String}\to\mathrm{String},\quad \mathbf{reverse}\colon\mathrm{String}\to\mathrm{String},\\
& \mathbf{equal}\colon\mathrm{Integer}\times\mathrm{Integer}\to\mathrm{Bool},\quad \mathbf{concat}\colon\mathrm{String}\times\mathrm{String}\to\mathrm{String},\\
& \mathbf{readString}\colon()\to\IO\,\mathrm{String},\quad \mathbf{putString}\colon\mathrm{String}\to\IO\,(),\\
& \mathbf{readFile}\colon\mathrm{String}\to\IO\,\mathrm{String},\quad \mathbf{split}\colon a\to a\times a,\\
& \mathbf{fst}\colon a\times a\to a,\quad \mathbf{twist}\colon a\times b\to b\times a,\\
& \mathbf{strength}\colon a\times\IO\,b\to\IO(a\times b),\quad \mathbf{pure}\colon a\to\IO\,a,
\end{align*}
napisać program $\mathbf{main}\colon()\to\IO\,()$, który:
\begin{enumerate}
\item wczytuje z klawiatury dwie nazwy plików, wczytuje pliki, sprawdza, czy ich zawartość ma tę samą długość, i wyświetla wynik na standardowym wyjściu;
\item wczytuje z klawiatury nazwę pliku, wczytuje plik i wyświetla jego nazwę zapisaną od tyłu złączoną z jego zawartością (zawartość wyświetlana jest nieodwrócona).
\end{enumerate}
\end{zadanie}

\section{Transformacje naturalne i monady}

\begin{zadanie}\label{zad:nat-ccc}
Niech $\C$ będzie kategorią kartezjańsko domkniętą. Pokazać, że następujące rodziny są transformacjami naturalnymi między odpowiednimi funktorami:
\begin{itemize}
\item $\{\eta_X\colon X\to(X\times A)^A\}_X$;
\item $\{\varepsilon_X\colon X^A\times A\to X\}_X$;
\item $\{\mu_X\colon ((X\times A)^A\times A)^A \to (X\times A)^A\}_X$, gdzie $\mu_X = (\varepsilon_{X\times A})^A$.
\end{itemize}
\end{zadanie}

\begin{zadanie}\label{zad:nat-set}
Pokazać, że następujące rodziny $\{\eta_X\}_X$ i $\{\mu_X\}_X$ są transformacjami naturalnymi między odpowiednimi $\Set$-funktorami:
\begin{enumerate}
\item $\{\eta_X\colon X\to\mathcal{P}X\}$, $\eta_X(x)=\{x\}$, oraz $\{\mu_X\colon\mathcal{P}\mathcal{P}X\to\mathcal{P}X\}$, $\mu_X(\mathcal{S}) = \bigcup\mathcal{S}$;
\item $\{\eta_X\colon X\to X^*\}$, $\eta_X(x)=x$, oraz $\{\mu_X\colon (X^*)^*\to X^*\}$, $\mu_X(s_1,\dots,s_n) = s_1 s_2 \dots s_n$;
\item dla ustalonego monoidu $M=(M,\cdot,1)$: $\{\eta_X\colon X\to M\times X\}$, $\eta_X(x)=(1,x)$, oraz $\{\mu_X\colon M\times M\times X\to M\times X\}$, $\mu_X(m,n,x)=(m\cdot n,\, x)$.
\end{enumerate}
\end{zadanie}

\begin{zadanie}
Pokazać, że następujące trójki są monadami:
\begin{enumerate}
\item $((\Id\times A)^A\colon\C\to\C,\, \mu,\, \eta)$, gdzie $\C$ jest CCC oraz $\mu,\eta$ są jak w Zadaniu~\ref{zad:nat-ccc};
\item $(T,\mu,\eta)$ dla $T\in\{\mathcal{P},\,(-)^*,\, M\times\Id\}$, gdzie $\mu,\eta$ są jak w Zadaniu~\ref{zad:nat-set}.
\end{enumerate}
\end{zadanie}

\begin{zadanie}
Niech $(T,\mu,\eta)$ będzie monadą w kategorii $\C$. Dla $f\colon A\to TB$, $g\colon B\to TC$ definiujemy
\[
g\cdot f\colon A\to TC, \qquad g\cdot f = \mu_C\circ T(g)\circ f.
\]
Pokazać, że dla $f,g$ oraz $h\colon C\to TD$ zachodzi
\[
(h\cdot g)\cdot f = h\cdot(g\cdot f), \qquad f\cdot\eta_A = \eta_B\cdot f = f.
\]
\end{zadanie}

\begin{zadanie}
Trójkę $(T,(-)^*,\eta)$ nazywamy \emph{trójką Kleisliego}, jeśli
\begin{itemize}
\item $T\colon\C_0\to\C_0$ jest przyporządkowaniem na obiektach;
\item $\eta = \{\eta_X\colon X\to TX\}_X$ jest rodziną strzałek;
\item $(-)^*$ przyporządkowuje każdej strzałce $f\colon X\to TY$ strzałkę $f^*\colon TX\to TY$,
\end{itemize}
spełniającymi równania
\[
\eta_X^* = \id_{TX}, \qquad f^*\circ\eta_X = f, \qquad (g^*\circ f)^* = g^*\circ f^*.
\]
Pokazać, że jeśli $(T,(-)^*,\eta)$ jest trójką Kleisliego, to $(T,\mu,\eta)$ jest monadą, gdzie $T(f) = (\eta_Y\circ f)^*$ dla $f\colon X\to Y$ oraz $\mu = \{\mu_X\colon T^2 X\to TX\}$ dla $\mu_X = (\id_{TX})^*$. Ponadto, jeśli $(T,\mu,\eta)$ jest monadą, to $(T,(-)^*,\eta)$ jest trójką Kleisliego, gdzie $f^* = \mu_Y\circ T(f)$ dla $f\colon X\to TY$.
\end{zadanie}


\end{document}
